% ---------------------------------------------------------------------------
% Chapitre 03 — Rappels d'interprétation abstraite
% Sources : notes/04-domains-lattices.md (§3.1, §3.3, §3.5, §4.3, §6.1, §10),
% notes/06-engine-fixpoint.md (§4.3, §4.4, §4.9, §4.12, §6.3, §6.4, §6.16),
% notes/09-rules-api.md (§3.3, §4.4), ADR-005, ADR-014, ADR-020, ADR-021, ADR-025.
% Sorties observées au commit e67b10a avec `cargo run -q -- check` sur les
% fichiers /tmp/ch03/*.tsx. Les valeurs abstraites intermédiaires viennent
% d'une sonde hors dépôt (/tmp/ch03/probe) : `probe` réplique la boucle externe
% de analyze_component_impl et appelle analyze_component ; `wl` réplique
% analyze_cfg en imprimant chaque visite de bloc ; `ops` évalue les
% opérations de domaine des exercices.
% ---------------------------------------------------------------------------
\chapter{Rappels d'interprétation abstraite}
\label{chap:interpretation-abstraite}

\begin{objectifs}
  \item Comprendre pourquoi un analyseur statique \emph{approxime}, et dans quel sens il a le droit de se tromper.
  \item Maîtriser le vocabulaire des treillis (ordre, join, meet, top, bottom, hauteur) et la concrétisation.
  \item Savoir calculer un point fixe par itération de Kleene et par worklist, et savoir pourquoi le widening est nécessaire.
  \item Distinguer un fait must d'un fait may, et relier cette polarité aux niveaux \Error{}, \Warning{}, \Info{}.
  \item Connaître la dominance, l'inlining, les mises à jour fortes et faibles, et situer \reactant{} parmi ces choix.
\end{objectifs}

\prerequis{\cref{chap:introduction} et \cref{chap:react-modele} (le cycle rendu → commit → effets, le batching des setters). Aucune connaissance préalable d'analyse statique n'est supposée ; la lecture de Rust est supposée.}

\section{Pourquoi approximer}
\label{sec:interpretation-abstraite-approximer}

\subsection{Une question que l'analyseur doit trancher}

Partons de deux composants presque identiques. Ils sont au cœur de ce
chapitre ; nous les suivrons jusqu'au bout.

\begin{lstlisting}[language=TypeScript,caption={Deux compteurs pilotés par un effet (\file{/tmp/ch03/counter.tsx}, repris de \file{tests/fixtures/widening.tsx})},label={lst:interpretation-abstraite-counter}]
import { useState, useEffect } from "react";

export function Runaway() {
  const [count, setCount] = useState(0);
  useEffect(() => {
    setCount(count + 1);
  }, [count]);
  return <div>{count}</div>;
}

export function Guarded() {
  const [count, setCount] = useState(0);
  useEffect(() => {
    if (count < 10) setCount(count + 1);
  }, [count]);
  return <div>{count}</div>;
}
\end{lstlisting}

À l'exécution (\cref{chap:react-modele}), \code{Runaway} ne s'arrête jamais :
chaque exécution de l'effet écrit \code{count + 1}, la dep \code{count} change,
l'effet repart. \code{Guarded} s'arrête après dix tours, quand la garde
\code{count < 10} devient fausse. La question que l'analyseur doit trancher
est donc : \emph{l'état \code{count} peut-il grandir sans borne ?}

Regardons ce que React fait vraiment, en notant l'ensemble des valeurs que
\code{count} prend au fil d'une exécution. Pour \code{Runaway}, c'est
$\{0,1,2,3,\dots\} = \Nat$ ; pour \code{Guarded}, c'est $\{0,1,\dots,10\}$. Ces
ensembles sont la réponse exacte. Si l'analyseur savait les calculer, il
n'aurait plus qu'à demander \og est-il fini ?\fg{}.

\begin{definition}{Sémantique concrète, sémantique collectrice}{semantique-concrete}
  La \terme{sémantique concrète} d'un programme décrit ses exécutions : suites
  d'états, un état donnant la valeur de chaque variable et de chaque slot
  d'état. Sa \terme{sémantique collectrice} en garde l'essentiel pour un
  analyste : pour chaque point du programme, l'\emph{ensemble} des états
  qu'une exécution quelconque peut y atteindre. On la note $\mathcal{C}$ ; un
  élément de $\mathcal{C}$ est une partie de l'ensemble des états.
\end{definition}

Pour \reactant{}, la sémantique concrète de référence est celle de React
décrite au \cref{chap:react-modele} (React-tRace étendue, \adr{1}) ; un
\og point\fg{} est un bloc du rendu, d'un effet ou d'un handler, et un état
comprend l'environnement local et la valeur de chaque slot.

\subsection{Pourquoi on ne peut pas calculer la réponse exacte}

Deux obstacles s'y opposent.

\begin{itemize}
  \item \textbf{La taille.} L'ensemble des valeurs de \code{count} dans
    \code{Runaway} est infini. On ne peut pas l'énumérer, et un ensemble quelconque
    d'entiers n'a pas de représentation finie.
  \item \textbf{L'indécidabilité.} Même en se limitant à des questions à
    réponse oui/non, la réponse exacte est hors d'atteinte. Le théorème de Rice
    affirme que toute propriété non triviale du comportement d'un programme
    (et \og ce composant boucle\fg{} en est une) est indécidable : aucun
    algorithme ne la tranche correctement sur tous les programmes. Un
    composant React contient un langage Turing-complet (JavaScript) ; on ne
    peut donc pas espérer un analyseur exact.
\end{itemize}

Il reste à renoncer à l'exactitude, mais \emph{dans un sens choisi}.
L'\terme{interprétation abstraite}~\cite{CousotCousot77,CousotCousot79} est la
théorie qui organise ce renoncement : on remplace les ensembles d'états par des
\emph{descriptions} finies, calculables, qui contiennent au moins tous les états
possibles. On calcule ces descriptions en \og exécutant\fg{} le programme sur
elles, et l'on raisonne sur le résultat. Le manuel de Rival et
Yi~\cite{RivalYi2020} en donne un exposé complet ; ce chapitre n'en retient que
ce dont \reactant{} se sert, et montre à chaque fois où le code l'implémente.

\subsection{Sur-approximer, et ce que cela coûte}

Choisir le \og sens\fg{} de l'erreur, c'est choisir entre deux familles
d'analyses (\cref{fig:interpretation-abstraite-approx}).

\begin{itemize}
  \item Une \terme{sur-approximation} calcule un ensemble $A$ qui
    \emph{contient} tous les comportements possibles $\mathcal{C}$ :
    $\mathcal{C} \subseteq A$. Si $A$ ne rencontre pas la zone des
    comportements fautifs, le programme est correct : le silence de l'analyse
    est une preuve. En revanche $A$ peut rencontrer la zone fautive alors que
    $\mathcal{C}$ ne la rencontre pas : l'analyse signale un défaut qui n'existe
    pas, c'est un \emph{faux positif}.
  \item Une \terme{sous-approximation} calcule un ensemble contenu dans
    $\mathcal{C}$ (c'est ce que fait un test : il exécute quelques
    comportements). Tout défaut qu'elle trouve est réel, mais son silence ne
    prouve rien : elle peut manquer un défaut, c'est un \emph{faux négatif}.
\end{itemize}

\begin{figure}[htbp]\centering
\begin{tikzpicture}[font=\small]
  % zone fautive
  \fill[ra-red!12] (3.1,-1.6) rectangle (6.2,1.6);
  \node[text=ra-red,font=\small\bfseries] at (5.35,1.25) {défaut};
  % sur-approximation A
  \draw[ra-amber,thick,dashed,fill=ra-amber!8] (1.2,0) ellipse (2.45 and 1.35);
  \node[text=ra-amber,anchor=south east] at (0.2,1.25) {$A$ (abstrait)};
  % sémantique collectrice C
  \draw[ra-blue,thick,fill=ra-blue!15] (0.9,-0.05) ellipse (1.5 and 0.8);
  \node[text=ra-blue] at (0.5,0) {$\mathcal{C}$ (concret)};
  % intersection A et défaut
  \node[etiquette,align=center] at (4.55,-0.25) {$A$ touche le défaut,\\$\mathcal{C}$ non :\\faux positif};
  \draw[fleche] (3.9,-0.9) -- (3.4,-0.55);
\end{tikzpicture}
\caption{Sur-approximation. Le concret $\mathcal{C}$ (ensemble des états
réellement atteignables) est inclus dans l'abstrait $A$. Tant que $A$ évite la
zone de défaut, le programme est prouvé correct ; si seul $A$ la touche,
l'analyse émet un faux positif. Jamais un défaut de $\mathcal{C}$ ne peut
échapper à $A$.}
\label{fig:interpretation-abstraite-approx}
\end{figure}

\reactant{} a fait le premier choix et l'a érigé en invariant (\file{CLAUDE.md}) :
\og faux positifs tolérés, faux négatifs INTERDITS\fg{}. Nous formaliserons cette
promesse à la \cref{sec:interpretation-abstraite-soundness}, une fois les outils
en place. Retenons dès maintenant qu'elle a une conséquence inattendue : quand
l'outil veut affirmer un défaut \emph{certain} (niveau \Error), il lui faut un
autre type de raisonnement, celui des faits \emph{must}
(\cref{sec:interpretation-abstraite-must-may}), car une sur-approximation ne
prouve jamais qu'un comportement a lieu.

\begin{remarque}
  Le mot \og approximation\fg{} ne veut pas dire \og à peu près\fg{}. Une
  sur-approximation est une \emph{borne} : elle est fausse au sens où elle
  contient des états impossibles, mais vraie au sens où elle contient tous les
  états possibles. Toute la suite du chapitre consiste à rendre cette borne
  calculable (treillis, points fixes, widening) sans jamais la rendre fausse
  dans l'autre sens.
\end{remarque}

\section{Ordres, treillis, top et bottom}
\label{sec:interpretation-abstraite-treillis}

\subsection{Un premier exemple : que sait-on d'un booléen ?}

Soit un composant qui lit une prop booléenne \code{open}. Sans rien savoir de
l'appelant, l'analyseur doit décrire les valeurs possibles de \code{open}. Il
y a quatre descriptions utiles :
\begin{itemize}
  \item \og \code{open} ne peut valoir que \code{true}\fg{} ;
  \item \og \code{open} ne peut valoir que \code{false}\fg{} ;
  \item \og \code{open} peut valoir l'un ou l'autre\fg{} : on ne sait rien, c'est
    l'information minimale, notée $\Top$ (\emph{top}) ;
  \item \og aucune valeur n'est possible\fg{} : le point est inatteignable,
    notée $\Bot$ (\emph{bottom}).
\end{itemize}
Ces descriptions sont ordonnées par la \emph{quantité de valeurs qu'elles
admettent} : $\Bot$ en admet le moins, $\Top$ le plus ; \code{true} et
\code{false} ne sont pas comparables. \Cref{fig:interpretation-abstraite-hasse-bool}
dessine cet ordre : un trait montant de $a$ vers $b$ signifie \og $a$ est plus
précis que $b$\fg{}. C'est un \terme{diagramme de Hasse}.

Si un chemin du rendu donne \code{open = true} et un autre \code{open = false},
ce que l'on sait après la jonction des deux chemins est la plus petite
description qui admette les deux : $\Top$. Cette opération de \og réunion
d'informations\fg{} s'appelle le \emph{join}.

\begin{figure}[htbp]\centering
\begin{tikzpicture}[treillis,node distance=9mm and 9mm]
  \node[label={[etiquette]right:$\gam = \emptyset$}] (bot) {$\Bot$};
  \node[above left=of bot,label={[etiquette]left:$\{\mathtt{true}\}$}] (t) {\code{True}};
  \node[above right=of bot,label={[etiquette]right:$\{\mathtt{false}\}$}] (f) {\code{False}};
  \node[above right=of t,label={[etiquette]right:$\{\mathtt{true},\mathtt{false}\}$}] (top) {$\Top$};
  \draw (bot) -- (t) -- (top) -- (f) -- (bot);
  \begin{scope}[xshift=7.2cm]
    \node (sbot) {$\Bot$};
    \node[above=13mm of sbot] (s2) {\code{One(c,1)}};
    \node[left=4mm of s2] (s1) {\code{One(c,0)}};
    \node[right=4mm of s2] (s3) {\code{One(d,0)}};
    \node[right=2mm of s3] (sdots) {$\cdots$};
    \node[above=13mm of s2] (stop) {$\Top$};
    \foreach \n in {s1,s2,s3} { \draw (sbot) -- (\n) -- (stop); }
  \end{scope}
\end{tikzpicture}
\caption{Deux treillis plats. À gauche, \code{BoolVal}
(\file{src/domains/impls/bool_val.rs}) avec la concrétisation de chaque élément
(\cref{sec:interpretation-abstraite-concretisation}) ; à droite, \code{SetterVal}
(\file{src/domains/impls/setter_val.rs}), dont les éléments intermédiaires
nomment un setter précis (composant, label de hook). Deux éléments
intermédiaires distincts sont incomparables ; leur join est $\Top$.}
\label{fig:interpretation-abstraite-hasse-bool}
\end{figure}

\subsection{Définitions}

\begin{definition}{Ordre partiel, treillis, join, meet}{treillis}
  Un \terme{ordre partiel} sur un ensemble $L$ est une relation $\lleq$
  réflexive, antisymétrique et transitive ; deux éléments peuvent être
  incomparables. Pour $a, b \in L$ :
  \begin{itemize}
    \item le \terme{join} $a \join b$ est la plus petite borne supérieure de
      $a$ et $b$ : $a \lleq a \join b$, $b \lleq a \join b$, et tout $c$ qui
      majore $a$ et $b$ vérifie $a \join b \lleq c$ ;
    \item le \terme{meet} $a \meet b$ est, dualement, la plus grande borne
      inférieure.
  \end{itemize}
  $(L, \lleq)$ est un \terme{treillis} si toute paire admet un join et un meet ;
  il est \emph{complet} si toute partie (même infinie) admet une borne
  supérieure (notée $\bigsqcup$) et une borne inférieure. Le plus petit élément, s'il
  existe, est le \terme{bottom} $\Bot$ ; le plus grand est le \terme{top} $\Top$.
  Une \terme{chaîne ascendante} est une suite $x_0 \lleq x_1 \lleq \cdots$ ; la
  \terme{hauteur} de $L$ est la longueur maximale d'une chaîne strictement
  croissante (elle peut être infinie).
\end{definition}

En analyse, on lit l'ordre comme un ordre d'\emph{information} : $a \lleq b$
signifie \og $a$ est plus précis que $b$\fg{}, ou \og $b$ admet au moins tous
les comportements de $a$\fg{}. $\Bot$ décrit un point inatteignable, $\Top$ un
point dont on ne sait rien (on dira aussi \og inconnu\fg{}). Le join est
l'opération des \emph{jonctions} du programme (après un \code{if}, en tête de
boucle, entre deux tours de React) : il fusionne ce que disent plusieurs
chemins. Le meet raffinerait une information par une autre ; nous verrons que
\reactant{} ne s'en sert presque pas.

\begin{exemple}{Le treillis plat des booléens}{treillis-plat}
  Dans \code{BoolVal}, $\code{True} \join \code{False} = \Top$,
  $\code{True} \meet \code{False} = \Bot$, $\Top \join \code{True} = \Top$ et
  $\Bot \lleq x \lleq \Top$ pour tout $x$. La hauteur est 2 (plus longue chaîne
  stricte : $\Bot \sqsubset \code{True} \sqsubset \Top$). Les valeurs ont été
  vérifiées sur le code (sonde \code{ops}, \code{True vs False = None},
  \code{True join False = Top}).
\end{exemple}

\subsection{Ce que fait le code : le trait \code{AbstractDomain}}

Dans \reactant{}, un domaine abstrait est un type Rust qui implémente le trait
\code{AbstractDomain}. Sa tête traduit exactement la
définition~\ref{def:treillis} :

\rustsnippet[label={lst:interpretation-abstraite-abstractdomain}]{src/domains/mod.rs}{16}{38}{le trait des domaines abstraits}

\begin{itemize}
  \item \textbf{L'ordre est \code{PartialOrd}.} Le supertrait
    \code{PartialOrd} porte l'ordre du treillis : \code{a <= b} signifie
    $a \lleq b$, et \code{partial_cmp} rend \code{None} pour deux éléments
    incomparables. Ce n'est pas un ordre total (\code{Ord} serait faux).
  \item \textbf{\code{bottom}, \code{top}, \code{join}, \code{meet}} sont les
    quatre ingrédients de la définition. \code{is_bottom} teste
    l'inatteignabilité (un domaine peut avoir plusieurs représentations de
    $\Bot$, c'est le cas des intervalles).
  \item \textbf{\code{widen} et \code{widen_to}} sont les opérateurs de
    widening de la \cref{sec:interpretation-abstraite-widening}. Par défaut
    \code{widen_to} ignore ses seuils et délègue à \code{widen}.
  \item \textbf{\code{PartialEq}} sert au test de convergence des points
    fixes.
\end{itemize}

\begin{piege}
  Le commentaire en tête du trait annonce \code{Clone + Copy}, mais la borne
  réelle est \code{Clone + PartialEq + PartialOrd + Debug} : pas de
  \code{Copy}. Le commentaire date d'avant le domaine produit (\adr{15}) et la
  stabilité versionnée (\adr{17}), dont les valeurs contiennent des ensembles.
  Se fier à la ligne \code{pub trait}, pas au commentaire.
\end{piege}

Le domaine le plus simple du code est le treillis plat des booléens :

\rustsnippet[label={lst:interpretation-abstraite-boolval}]{src/domains/impls/bool_val.rs}{1}{13}{\code{BoolVal}, un treillis plat}

La macro \code{flat_lattice!} engendre l'ordre et les opérations de tout
treillis plat (un $\Bot$, un $\Top$, des éléments intermédiaires deux à deux
incomparables). Son \code{partial_cmp} et son \code{join} sont la
définition~\ref{def:treillis} écrite en \code{match} :

\rustsnippet[label={lst:interpretation-abstraite-flat}]{src/domains/impls/mod.rs}{10}{40}{la macro \code{flat_lattice!} (ordre et join)}

L'ordre des bras compte : l'égalité est testée d'abord, puis
\og $\Bot$ à gauche ou $\Top$ à droite\fg{} (donc \code{Less}) ; le cas
$(\Bot, \Top)$ tombe ainsi sur \code{Less}. Deux éléments intermédiaires
différents n'ont aucun bras : \code{None}. Le join suit le même schéma : égaux,
on garde ; l'un est $\Bot$, on garde l'autre ; sinon $\Top$. Le \code{meet}
(\src{src/domains/impls/mod.rs}{41}{47}) est le dual, et \code{widen} est le join
(\src{src/domains/impls/mod.rs}{48}{50}) : un treillis de hauteur finie n'a pas
besoin de mieux (\cref{sec:interpretation-abstraite-widening}).

Les stores comparent leurs valeurs slot par slot avec un seul utilitaire, qui
dit bien que l'ordre est partiel :

\rustsnippet[label={lst:interpretation-abstraite-leq-pointwise}]{src/domains/stores/mod.rs}{19}{25}{$a \lleq b$ dans un ordre partiel}

\begin{piege}
  Écrire \code{!(a > b)} pour tester $a \lleq b$ est faux dans un ordre
  partiel : deux éléments incomparables ne sont ni \code{>} ni \code{<=}.
  \code{leq_pointwise} ne retient que \code{Less} ou \code{Equal}, et un
  couple incomparable n'est donc \emph{pas} considéré comme convergé.
\end{piege}

\subsection{Autres treillis du projet, en un coup d'œil}

Le \cref{chap:treillis-simples} détaillera les domaines de base ; il suffit ici
d'en connaître la forme.
\begin{itemize}
  \item \code{SetterVal} : treillis plat dont les éléments intermédiaires sont
    les setters \code{One(ComponentId, HookLabel)}
    (\cref{fig:interpretation-abstraite-hasse-bool}, à droite).
  \item \code{StrConst} : ensembles d'au plus quatre chaînes constantes, ordonnés
    par inclusion, avec $\Top$ au-delà ; hauteur finie.
  \item \code{Interval} : intervalles de nombres, ordonnés par inclusion ;
    hauteur \emph{infinie} (\cref{sec:interpretation-abstraite-compteur}).
  \item \code{Stability} : la stabilité référentielle versionnée d'\adr{17}
    (\cref{chap:statevalue-stabilite}).
  \item \code{StateValue} : le produit des précédents, une composante par sorte
    de valeur JavaScript (\adr{15}). Un produit de treillis est un treillis,
    ordonné composante par composante.
\end{itemize}

\section{Concrétisation et correction}
\label{sec:interpretation-abstraite-concretisation}

\subsection{Ce que \og veut dire\fg{} une valeur abstraite}

Dans \code{Guarded}, l'analyseur finit par attribuer à \code{count}
l'intervalle $[0,10]$. Pour que cette phrase soit une affirmation vérifiable,
il faut dire \emph{quels nombres} elle autorise : les entiers de $0$ à $10$. La
fonction qui associe à chaque valeur abstraite l'ensemble des valeurs concrètes
qu'elle décrit est la concrétisation.

\begin{definition}{Concrétisation, abstraction, correction}{concretisation}
  Soit $(L, \lleq)$ un domaine abstrait et $\mathcal{D}$ l'ensemble des
  valeurs (ou des états) concrets. La \terme{concrétisation}
  $\gam : L \to \powerset(\mathcal{D})$ associe à chaque valeur abstraite
  l'ensemble des valeurs concrètes qu'elle représente. Elle est
  \emph{monotone} : $a \lleq b \Rightarrow \gam(a) \subseteq \gam(b)$. On
  demande $\gam(\Bot) = \emptyset$ et, souvent, $\gam(\Top) = \mathcal{D}$.

  Une \terme{abstraction} $\abs : \powerset(\mathcal{D}) \to L$ associe à un
  ensemble concret la meilleure description disponible. Quand
  $\abs(X) \lleq a \iff X \subseteq \gam(a)$, le couple $(\abs, \gam)$ est une
  \terme{correspondance de Galois}.

  Une opération abstraite $f^\sharp : L \to L$ est \terme{correcte} vis-à-vis
  d'une opération concrète $f$ si, pour tout $a$,
  $f(\gam(a)) \subseteq \gam(f^\sharp(a))$ : ce que $f$ peut produire à partir
  de valeurs décrites par $a$ est décrit par $f^\sharp(a)$. Pour le join :
  $\gam(a) \cup \gam(b) \subseteq \gam(a \join b)$.
\end{definition}

La condition de correction est la seule qui compte pour la soundness : une
opération abstraite peut être imprécise (rendre plus que nécessaire), jamais
rendre moins. Une abstraction $\abs$ n'est qu'une commodité de présentation :
\reactant{} ne définit aucune fonction $\abs$ générale. Il abstrait les
littéraux au cas par cas (\code{Interval::point(0.0)} pour \code{0}) et calcule
tout le reste par des opérations abstraites ; $\gam$ n'apparaît, lui, que dans
les commentaires du code, sous forme de phrases.

\begin{exemple}{Deux concrétisations}{gamma}
  \begin{itemize}
    \item \code{BoolVal} : $\gam(\Bot) = \emptyset$,
      $\gam(\code{True}) = \{\mathtt{true}\}$,
      $\gam(\code{False}) = \{\mathtt{false}\}$,
      $\gam(\Top) = \{\mathtt{true}, \mathtt{false}\}$. Ici $\gam$ est une
      bijection avec les parties de $\{\mathtt{true},\mathtt{false}\}$ ; le
      join est \emph{exact} : $\gam(a \join b) = \gam(a) \cup \gam(b)$.
    \item \code{Interval} : $\gam([a,b]) = \{x \mid a \le x \le b\}$, restreint
      aux entiers si l'intervalle est marqué entier. Le join
      $[0,0] \join [10,10] = [0,10]$ est correct
      ($\{0\} \cup \{10\} \subseteq \gam([0,10])$) mais \emph{pas} exact :
      $\gam([0,10])$ contient aussi $1, \dots, 9$. C'est la perte de précision
      inhérente à un domaine convexe.
  \end{itemize}
\end{exemple}

\subsection{Ce que fait le code : l'intervalle}

La structure \code{Interval} documente sa propre concrétisation :

\rustsnippet[label={lst:interpretation-abstraite-interval}]{src/domains/impls/interval.rs}{5}{27}{les intervalles et leur concrétisation}

\begin{itemize}
  \item \code{lo}, \code{hi} : bornes fermées, en \code{f64} ; $\pm\infty$
    représente l'absence de borne. Un intervalle avec \code{lo > hi} est
    $\Bot$ (vide).
  \item \code{is_int} : \og toute valeur concrète est entière\fg{}. C'est une
    annotation de \emph{précision} : elle autorise le raffinement
    $x < 5 \Rightarrow x \le 4$, que nous retrouverons avec le widening. Elle
    est volontairement exclue de l'égalité (et de l'ordre) : le point fixe
    converge quelles que soient les valeurs de ce bit.
\end{itemize}

Le join des intervalles est l'enveloppe convexe (\emph{hull}) :

\rustsnippet[label={lst:interpretation-abstraite-hull}]{src/domains/impls/interval.rs}{69}{83}{le join des intervalles}

Sa correction se lit ligne à ligne : si $x \in \gam(a)$, alors
$\min(a.lo, b.lo) \le a.lo \le x \le a.hi \le \max(a.hi, b.hi)$ ; et si les deux
opérandes ne contiennent que des entiers, leur réunion aussi. Les deux premiers
tests traitent $\Bot$ comme élément neutre, ce que demande
$\gam(\Bot) = \emptyset$.

\begin{piege}
  $\Bot$ a plusieurs représentations. Le $\Bot$ canonique est
  $[+\infty, -\infty]$, mais un raffinement peut produire, par exemple,
  $[5, 0]$ (c'est ce que la sonde \code{wl} observe dans la boucle de la
  \cref{sec:interpretation-abstraite-point-fixe}). Les combinateurs testent
  \code{is_bottom()} d'abord et restent corrects ; en revanche
  \code{PartialEq} et \code{PartialOrd} comparent les bornes brutes, et
  $[5,0] \neq [5,1]$ pour eux alors que les deux sont vides. Conséquence
  pratique : une worklist peut revisiter un bloc inatteignable pour rien.
  C'est une imprécision de performance, pas une faute de soundness.
\end{piege}

\subsection{\og Absent\fg{} vaut $\Top$ ou $\Bot$ selon le contexte}

Deux conteneurs de \reactant{} donnent une valeur par défaut aux clefs
absentes, et ils ne donnent pas la même.

\rustsnippet[label={lst:interpretation-abstraite-lookup}]{src/domains/stores/abstract_env.rs}{78}{81}{une variable inconnue vaut $\Top$}

\rustsnippet[label={lst:interpretation-abstraite-store-get}]{src/domains/stores/state_store.rs}{24}{27}{un slot jamais écrit vaut $\Bot$}

Les deux choix sont corrects, parce que les deux questions diffèrent. Une
variable que l'environnement ne connaît pas (un import non résolu, un global
du navigateur) peut valoir n'importe quoi : répondre autre chose que $\Top$
inventerait une information. Un slot d'état, au contraire, ne reçoit de valeur
que par ses écritures ; le store accumule ces écritures par join à partir de
$\Bot$, et le moteur l'amorce dès le départ avec la valeur initiale de chaque
\code{useState} (\cref{chap:point-fixe-composant}). \og Pas d'écriture
vue\fg{} y signifie \og pas de valeur ajoutée\fg{}, donc $\Bot$.

\begin{piege}
  Confondre les deux conventions est une source classique de faux négatif :
  lire un slot absent comme $\Bot$ \emph{dans une règle} reviendrait à dire
  \og cet état ne prend aucune valeur\fg{}, donc \og il ne change pas\fg{}. Les
  règles lisent l'état convergé, où tout slot a été amorcé ; hors de ce cadre,
  une absence est une absence d'information, et l'information minimale est
  $\Top$.
\end{piege}

\section{Exemple filé : le compteur en intervalles}
\label{sec:interpretation-abstraite-compteur}

\subsection{Le treillis des intervalles}

Pour répondre à la question de la
\cref{sec:interpretation-abstraite-approximer}, un domaine qui décrit des
ensembles de nombres par leurs bornes suffit : le treillis des intervalles
(\cref{fig:interpretation-abstraite-hasse-intervalles}). Son ordre est
l'inclusion, $[a,b] \lleq [c,d] \iff c \le a \wedge b \le d$
(\src{src/domains/impls/interval.rs}{312}{328}) ; son join est l'enveloppe
convexe, son meet l'intersection. Deux intervalles qui se chevauchent sans
s'inclure, comme $[1,3]$ et $[2,9]$, sont incomparables (\code{partial_cmp}
rend \code{None}, vérifié avec la sonde \code{ops}).

\begin{figure}[htbp]\centering
\begin{tikzpicture}[treillis,x=13mm,y=11mm]
  \node (bot) at (1,0) {$\Bot$};
  \node (p0) at (0,1) {$[0,0]$};
  \node (p1) at (1,1) {$[1,1]$};
  \node (p2) at (2,1) {$[2,2]$};
  \node (pd) at (3,1) {$\cdots$};
  \node (a01) at (0.5,2) {$[0,1]$};
  \node (a12) at (1.5,2) {$[1,2]$};
  \node (a02) at (1,3) {$[0,2]$};
  \node (dots) at (1,4) {$\vdots$};
  \node (a0i) at (1,5) {$[0,+\infty]$};
  \node (ami) at (-0.6,5) {$[-\infty,0]$};
  \node (top) at (0.2,6) {$\Top = [-\infty,+\infty]$};
  \draw (bot) -- (p0); \draw (bot) -- (p1); \draw (bot) -- (p2);
  \draw (p0) -- (a01); \draw (p1) -- (a01); \draw (p1) -- (a12); \draw (p2) -- (a12);
  \draw (a01) -- (a02); \draw (a12) -- (a02);
  \draw (a02) -- (dots); \draw (dots) -- (a0i);
  \draw (a0i) -- (top); \draw (ami) -- (top);
  \draw[dotted] (p0) -- (ami);
  \draw[ra-blue,very thick,->] (2.3,1.2) -- node[right,etiquette,align=left]{chaîne ascendante\\infinie $[0,0] \lleq [0,1] \lleq [0,2] \lleq \cdots$} (2.3,4.8);
\end{tikzpicture}
\caption{Le treillis des intervalles, tronqué. Un intervalle est au-dessus de
ceux qu'il contient. La chaîne $[0,0] \lleq [0,1] \lleq [0,2] \lleq \cdots$ ne
s'arrête jamais : la hauteur est infinie, ce qui obligera au widening
(\cref{sec:interpretation-abstraite-widening}).}
\label{fig:interpretation-abstraite-hasse-intervalles}
\end{figure}

Les opérations arithmétiques se relèvent borne à borne. Pour l'addition,
$[a,b] + [c,d] = [a+c, b+d]$ : si $x \in [a,b]$ et $y \in [c,d]$, alors
$a + c \le x + y \le b + d$. Le code est la formule
(\src{src/domains/impls/interval.rs}{148}{157}), précédée du cas $\Bot$ : une
somme dont un opérande est inatteignable est inatteignable.

\subsection{Un tour abstrait de \code{Runaway}}

Exécutons \og à la main\fg{} un tour de React sur des intervalles, comme le
fait le moteur (\cref{chap:point-fixe-composant} donnera le détail).
\begin{enumerate}
  \item \textbf{Amorçage.} \code{useState(0)} : le slot 0 (\code{count}) vaut
    $[0,0]$.
  \item \textbf{Rendu.} La variable \code{count} lit le slot : $[0,0]$. Le
    rendu n'écrit aucun état.
  \item \textbf{Effet.} Le corps évalue \code{count + 1}, soit
    $[0,0] + [1,1] = [1,1]$, et appelle le setter.
  \item \textbf{Écriture.} Le setter \emph{ne remplace pas} la valeur du slot :
    il la joint, $[0,0] \join [1,1] = [0,1]$. Nous justifierons ce choix à la
    \cref{sec:interpretation-abstraite-mises-a-jour} ; retenons que le slot
    décrit toutes les valeurs que \code{count} a pu prendre jusqu'ici.
\end{enumerate}
Au tour suivant, le rendu lit $[0,1]$, l'effet écrit $[1,2]$, le slot devient
$[0,2]$, et ainsi de suite. Au tour $k$ (en partant de 0), le slot vaut
$[0, k+1]$ : c'est exactement l'ensemble des valeurs que React donne à
\code{count} au cours des $k+1$ premières exécutions de l'effet. Chaque tour
abstrait est correct (il contient ce que React peut faire), mais la suite ne
se stabilise jamais.

Pour \code{Guarded}, un ingrédient s'ajoute : la garde. Sur la branche
\og vraie\fg{} de \code{count < 10}, l'analyseur sait que \code{count} est
strictement inférieur à 10 ; sur un intervalle d'entiers, il le
\emph{raffine} en $[0,k] \meet [-\infty, 9]$ (le code calcule ce meet
directement, par \code{Interval::narrow_lt}). C'est le \terme{raffinement de
garde} (\code{narrow_env_for_branch}, \cref{chap:cfg-analyzer-dominance}). Tant
que $k < 10$, le raffinement ne change rien. Au tour où le slot vaut
$[0,10]$, la branche vraie ne voit plus que $[0,9]$, l'effet écrit $[1,10]$, et
$[0,10] \join [1,10] = [0,10]$ : le tour n'ajoute plus rien.

\begin{table}[htbp]\centering\small
\begin{tabular}{@{}rllll@{}}
  \toprule
  tour & \multicolumn{2}{l}{\code{Runaway}} & \multicolumn{2}{l}{\code{Guarded}} \\
  \cmidrule(lr){2-3}\cmidrule(lr){4-5}
       & slot lu & slot après le tour & slot lu & slot après le tour \\
  \midrule
  0 & $[0,0]$ & $[0,1]$ & $[0,0]$ & $[0,1]$ \\
  1 & $[0,1]$ & $[0,2]$ & $[0,1]$ & $[0,2]$ \\
  2 & $[0,2]$ & $[0,3]$ & $[0,2]$ & $[0,3]$ \\
  $\vdots$ & & & & \\
  9 & $[0,9]$ & $[0,10]$ & $[0,9]$ & $[0,10]$ \\
  10 & $[0,10]$ & $[0,11]$ & $[0,10]$ & $[0,10]$ \quad(stable) \\
  $\vdots$ & & & & \\
  29 & $[0,29]$ & $[0,30]$ & & \\
  \bottomrule
\end{tabular}
\caption{Les tours abstraits des deux compteurs \emph{sans widening}. Valeurs
observées avec la sonde \code{probe}, qui réplique la boucle externe du moteur
(\src{src/engine/fixpoint.rs}{355}{530}) en désactivant le widening
(\code{WT=1000}) : \code{Guarded} s'arrête au tour 10, \code{Runaway} croît
encore au tour 29, dernier tour exécuté par la sonde.}
\label{tab:interpretation-abstraite-tours}
\end{table}

\Cref{tab:interpretation-abstraite-tours} pose les deux questions des sections
suivantes. D'abord : quand la suite \emph{se stabilise}, comme pour
\code{Guarded}, pourquoi l'intervalle obtenu est-il une réponse correcte ? C'est
la théorie du point fixe (\cref{sec:interpretation-abstraite-point-fixe}).
Ensuite : quand elle \emph{ne se stabilise pas}, comme pour \code{Runaway}, comment
s'arrêter en temps fini sans perdre la correction ? C'est le widening
(\cref{sec:interpretation-abstraite-widening}).

\section{Transfert, équations de flot, point fixe}
\label{sec:interpretation-abstraite-point-fixe}

\subsection{Une boucle locale}

Le compteur de la section précédente itère sur les \emph{tours de React}.
Une boucle écrite dans le code pose le même problème à plus petite échelle,
à l'intérieur d'un seul corps de fonction. Le programme suivant est repris de
la fixture \file{tests/fixtures/widening.tsx} :

\begin{lstlisting}[language=TypeScript,caption={Une boucle bornée dans un effet (\file{/tmp/ch03/loop.tsx}, identique à \code{BoundedLocalLoop} de \file{tests/fixtures/widening.tsx})},label={lst:interpretation-abstraite-loop}]
import { useState, useEffect } from "react";

export function BoundedLocalLoop() {
  const [total, setTotal] = useState(0);
  useEffect(() => {
    let i = 0;
    while (i < 5) {
      i = i + 1;
    }
    setTotal(i);
  }, []);
  return <div>{total}</div>;
}
\end{lstlisting}

Concrètement, \code{i} vaut $0, 1, \dots, 5$ en tête de boucle et $5$ à la
sortie ; l'effet écrit $5$. L'analyseur doit retrouver au moins cela, sans
exécuter la boucle un nombre de fois qu'il ne connaît pas à l'avance. Le
\cref{chap:lowering-cfg} explique comment le corps de l'effet devient un
graphe de flot de contrôle (CFG, \cref{def:cfg}) ;
\cref{fig:interpretation-abstraite-cfg-loop} en donne la forme.

\begin{figure}[htbp]\centering
\begin{tikzpicture}[node distance=7mm and 16mm]
  \node[cfgbloc] (b0) {$B_0$ (entrée)\\let i = 0};
  \node[cfgbloc,below=of b0] (b1) {$B_1$ (en-tête)\\branch i < 5};
  \node[cfgbloc,below left=of b1] (b2) {$B_2$ (corps)\\i = i + 1};
  \node[cfgbloc,below right=of b1] (b3) {$B_3$ (sortie)\\setTotal(i)\\return};
  \draw[fleche] (b0) -- (b1);
  \draw[fleche] (b1) -- node[etiquette,left]{vrai} (b2);
  \draw[fleche] (b1) -- node[etiquette,right]{faux} (b3);
  \draw[flechemay] (b2.west) to[out=150,in=180] node[etiquette,left]{\code{Back}} (b1.west);
\end{tikzpicture}
\caption{CFG schématique du corps de l'effet de
\cref{lst:interpretation-abstraite-loop}. L'arc en pointillés est l'arc arrière
de la boucle ; le lowering le marque \code{EdgeKind::Back}
(\cref{chap:lowering-cfg}). La numérotation des blocs est celle de l'exposé,
pas nécessairement celle du lowering.}
\label{fig:interpretation-abstraite-cfg-loop}
\end{figure}

\subsection{Fonctions de transfert}

À chaque bloc, l'analyseur associe une \terme{fonction de transfert}
abstraite : elle prend la valeur abstraite à l'entrée du bloc (ici, un
environnement qui donne un intervalle à \code{i}) et rend celle de la
sortie. Pour $B_2$, \code{i = i + 1} devient
$f_2^\sharp(\rho) = \rho[\code{i} \mapsto \rho(\code{i}) + [1,1]]$, avec
l'addition d'intervalles de la \cref{sec:interpretation-abstraite-compteur}.
Pour une arête issue d'une garde, la fonction de transfert est le
raffinement : sur l'arête \og vrai\fg{} de $B_1$, $\rho(\code{i})$ devient
$\rho(\code{i}) \meet [-\infty, 4]$ (\code{i} est entier), sur l'arête
\og faux\fg{}, $\rho(\code{i}) \meet [5, +\infty]$.

\begin{definition}{Fonction monotone, fonction correcte}{transfert}
  Une fonction $f^\sharp : L \to L$ est \terme{monotone} si
  $a \lleq b \Rightarrow f^\sharp(a) \lleq f^\sharp(b)$ : plus d'information en
  entrée ne donne pas moins d'information en sortie. Elle est
  \emph{correcte} vis-à-vis de l'instruction concrète $f$ au sens de la
  \cref{def:concretisation} : $f(\gam(a)) \subseteq \gam(f^\sharp(a))$.
\end{definition}

L'addition d'intervalles est monotone et correcte (la preuve tient en une
ligne, donnée plus haut). Le raffinement de garde est correct parce qu'il ne
retire que des valeurs pour lesquelles la garde est fausse. La monotonie
est la propriété confortable des manuels ; la correction est la seule dont
dépend la soundness, et c'est elle que les commentaires du code revendiquent
(\og \emph{Sound: result ⊒ self.hull(other)}\fg{} sur
\code{widen_to}, \src{src/domains/impls/interval.rs}{108}{113}).

\subsection{Équations de flot}

Posons une inconnue $X_i$ par bloc, la valeur abstraite à son entrée. Le CFG
donne directement un système d'équations :
\[
\begin{aligned}
  X_0 &= \rho_0 \quad \text{(l'environnement d'entrée de l'effet)}\\
  X_1 &= f_0^\sharp(X_0) \;\join\; f_2^\sharp(X_2)\\
  X_2 &= \mathit{vrai}^\sharp_{\code{i<5}}\big(X_1\big)\\
  X_3 &= \mathit{faux}^\sharp_{\code{i<5}}\big(X_1\big)
\end{aligned}
\]
L'en-tête $B_1$ a deux prédécesseurs : l'entrée et l'arc arrière. Son
équation les \emph{joint}. Regroupons les inconnues en un vecteur
$X = (X_0, \dots, X_3)$ et les membres droits en une fonction $F^\sharp$ ; le
système s'écrit $X = F^\sharp(X)$. Une solution est un \emph{point fixe} de
$F^\sharp$.

\begin{definition}{Point fixe, post-point fixe, itération de Kleene}{point-fixe}
  Soit $F : L \to L$ sur un treillis $L$. Un \terme{point fixe} est un $x$ tel
  que $F(x) = x$ ; le \terme{plus petit point fixe} se note $\lfp F$. Un
  \terme{post-point fixe} est un $x$ tel que $F(x) \lleq x$ : appliquer $F$
  n'ajoute plus rien. L'\terme{itération de Kleene} est la suite
  \[ x_0 = \Bot, \qquad x_{n+1} = F(x_n) \;(\text{ou, en accumulant, } x_n \join F(x_n)). \]
\end{definition}

\begin{theoreme}{Tarski, Kleene}{kleene}
  Si $L$ est un treillis complet et $F$ monotone, alors $F$ a un plus petit
  point fixe, et c'est le plus petit post-point fixe :
  $\lfp F$ est la borne inférieure de $\{ x \mid F(x) \lleq x \}$ (Tarski). Si de plus $F$
  préserve les bornes supérieures des chaînes, la suite de Kleene est
  croissante et $\lfp F = \bigsqcup_{n} x_n$ ; si $L$ est de hauteur finie,
  elle atteint $\lfp F$ en un nombre fini de pas.
\end{theoreme}

La conséquence qui compte pour un analyseur est la suivante.

\begin{theoreme}{Un post-point fixe abstrait est sound}{post-point-fixe}
  Soit $F$ la fonction concrète de la sémantique collectrice (sur les parties
  d'états) et $F^\sharp$ une fonction abstraite correcte :
  $F(\gam(a)) \subseteq \gam(F^\sharp(a))$ pour tout $a$. Si $a$ est un
  post-point fixe de $F^\sharp$, alors $\lfp F \subseteq \gam(a)$.
\end{theoreme}

\noindent\emph{Preuve.} Par correction puis monotonie de $\gam$,
$F(\gam(a)) \subseteq \gam(F^\sharp(a)) \subseteq \gam(a)$ : $\gam(a)$ est un
post-point fixe de $F$, donc il contient le plus petit (Tarski). \qed

Ce théorème est précieux parce qu'il ne demande \emph{ni} que $F^\sharp$ soit
monotone, \emph{ni} que $a$ soit le plus petit point fixe abstrait. Il suffit
de s'arrêter sur un $a$ tel que $F^\sharp(a) \lleq a$ : on peut y arriver par
n'importe quel chemin, y compris en sautant des étapes (c'est ce que fera le
widening). Tout le moteur de \reactant{} repose sur ce seul test ;
\cref{alg:interpretation-abstraite-kleene} en donne la forme la plus simple.

\begin{algorithm}[htbp]
\KwIn{$F^\sharp : L \to L$}
\KwOut{un post-point fixe $x$ de $F^\sharp$}
$x \gets \Bot$\;
\Repeat{$y \lleq x$}{
  $y \gets F^\sharp(x)$\;
  \lIf{$y \lleq x$}{\Return{$x$}}
  $x \gets x \join y$\;
}
\caption{Itération de Kleene (forme accumulante)}
\label{alg:interpretation-abstraite-kleene}
\end{algorithm}

\subsection{La worklist : n'itérer que ce qui change}

Itérer $F^\sharp$ en bloc recalcule tous les blocs à chaque pas, même ceux
dont l'entrée n'a pas bougé. Une \terme{worklist} ne recalcule qu'un bloc dont
l'entrée a changé, et ne propage qu'aux successeurs concernés
(\cref{alg:interpretation-abstraite-worklist}). C'est une stratégie
d'itération chaotique : elle calcule le même résultat que Kleene quand $F^\sharp$
est monotone et le domaine de hauteur finie, avec beaucoup moins de travail.

\begin{algorithm}[htbp]
\KwIn{un CFG, un environnement d'entrée $\rho_0$, un seuil $k$}
\KwOut{l'environnement de sortie de chaque bloc atteint}
$\mathit{In}[\mathit{entry}] \gets \rho_0$ ; $W \gets [\mathit{entry}]$\;
\While{$W$ non vide}{
  $b \gets$ tête de $W$ (FIFO)\;
  $\mathit{Out}[b] \gets f_b^\sharp(\mathit{In}[b])$\;
  \ForEach{successeur $s$ de $b$, avec l'environnement raffiné $\rho_s$}{
    \uIf{$\mathit{In}[s]$ absent}{$\nu \gets \rho_s$\;}
    \uElseIf{$b \to s$ est un arc arrière et c'est son $k$-ième passage ou plus}{$\nu \gets \mathit{In}[s] \widen \rho_s$\;}
    \Else{$\nu \gets \mathit{In}[s] \join \rho_s$\;}
    \lIf{$\nu \neq \mathit{In}[s]$}{$\mathit{In}[s] \gets \nu$ ; enfiler $s$ s'il n'est pas déjà dans $W$}
  }
}
\caption{Worklist d'un CFG, telle que l'implémente \code{analyze_cfg}}
\label{alg:interpretation-abstraite-worklist}
\end{algorithm}

\subsection{Ce que fait le code : \code{analyze_cfg}}

\code{analyze_cfg} (\file{src/engine/cfg_analyzer.rs}) est la
\cref{alg:interpretation-abstraite-worklist} en Rust. Le
\cref{chap:cfg-analyzer-dominance} la détaille ; en voici les trois temps.
D'abord l'initialisation :

\rustsnippet[label={lst:interpretation-abstraite-acfg-init}]{src/engine/cfg_analyzer.rs}{47}{58}{\code{analyze_cfg} : initialisation}

\begin{itemize}
  \item \code{entry_envs} est la table $\mathit{In}$, \code{exit_envs} la table
    $\mathit{Out}$. Seul le bloc d'entrée reçoit un environnement : un bloc que
    rien n'atteint n'est jamais visité et n'a pas d'entrée dans
    \code{exit_envs}, ce qui revient à lui donner $\Bot$.
  \item \code{state_out} est \emph{une seule} copie du store d'état pour toute
    la passe. Nous y reviendrons à la
    \cref{sec:interpretation-abstraite-reactant} : l'état n'est pas suivi bloc
    par bloc.
  \item \code{in_worklist} évite d'enfiler deux fois le même bloc ;
    \code{back_edge_counts} compte les passages par arc arrière, bloc cible par
    bloc cible (ce compteur déclenchera le widening).
\end{itemize}

Ensuite l'application de la fonction de transfert d'un bloc :

\rustsnippet[label={lst:interpretation-abstraite-acfg-body}]{src/engine/cfg_analyzer.rs}{60}{89}{\code{analyze_cfg} : transfert d'un bloc}

Les instructions du bloc sont exécutées une à une par le \code{Transfer}
(\cref{chap:transfert-stores}) ; si le bloc se termine par \code{return e}, les
effets de bord de \code{e} sont aussi exécutés (un effet
\code{() => setN(1)} a son appel dans le \code{return}). Enfin, la
propagation :

\rustsnippet[label={lst:interpretation-abstraite-acfg-succ}]{src/engine/cfg_analyzer.rs}{91}{121}{\code{analyze_cfg} : propagation aux successeurs}

\code{outgoing} (\src{src/engine/cfg_analyzer.rs}{126}{150}) rend chaque
successeur avec l'environnement raffiné par la garde du \code{Branch}. La
nouvelle entrée est l'environnement propagé s'il n'y en avait pas, son join
avec l'ancienne sinon, et son widening à seuils sur un arc arrière dès le
\code{widen_threshold}-ième passage (3 par défaut,
\src{src/engine/fixpoint.rs}{51}{60}). Le successeur n'est ré-enfilé que si son
entrée a changé.

\begin{exemple}{La worklist sur la boucle locale}{worklist-boucle}
  Déroulons \code{analyze_cfg} sur \cref{fig:interpretation-abstraite-cfg-loop},
  en ne montrant que la composante numérique de \code{i}. Les seuils du
  composant sont les littéraux numériques du programme, $T = \{0, 1, 5\}$
  (\cref{sec:interpretation-abstraite-widening}). Le déroulé est fait à la
  main d'après le code ; son résultat final est épinglé par le test
  \code{bounded_local_loop_is_precise} (\src{tests/widening_e2e.rs}{124}{137}).

  \smallskip
  \begin{tabular}{@{}rlll@{}}
    \toprule
    pas & bloc & entrée → sortie & effet sur les successeurs \\
    \midrule
    1 & $B_0$ & $\rho_0$ → $[0,0]$ & $\mathit{In}[B_1] := [0,0]$ \\
    2 & $B_1$ & $[0,0]$ & $\mathit{In}[B_2] := [0,0]$, $\mathit{In}[B_3] := [5,0]$ ($\Bot$) \\
    3 & $B_2$ & $[0,0]$ → $[1,1]$ & arrière n\textsuperscript{o}1 : $[0,0] \join [1,1] = [0,1]$ \\
    4 & $B_3$ & $\Bot$ & \code{setTotal} n'ajoute rien \\
    5 & $B_1$ & $[0,1]$ & $\mathit{In}[B_2] := [0,1]$, $\mathit{In}[B_3] := [5,1]$ ($\Bot$) \\
    6 & $B_2$ & $[0,1]$ → $[1,2]$ & arrière n\textsuperscript{o}2 : $[0,1] \join [1,2] = [0,2]$ \\
    7 & $B_3$ & $\Bot$ & rien \\
    8 & $B_1$ & $[0,2]$ & $\mathit{In}[B_2] := [0,2]$, $\mathit{In}[B_3] := [5,2]$ ($\Bot$) \\
    9 & $B_2$ & $[0,2]$ → $[1,3]$ & arrière n\textsuperscript{o}3 : $[0,2] \widen_T [1,3] = [0,5]$ \\
    10 & $B_3$ & $\Bot$ & rien \\
    11 & $B_1$ & $[0,5]$ & $\mathit{In}[B_2] := [0,4]$, $\mathit{In}[B_3] := [5,5]$ \\
    12 & $B_2$ & $[0,4]$ → $[1,5]$ & arrière n\textsuperscript{o}4 : $[0,5] \widen_T [1,5] = [0,5]$, inchangé \\
    13 & $B_3$ & $[5,5]$ & \code{setTotal} écrit $[5,5]$ ; le slot devient $[0,0] \join [5,5] = [0,5]$ \\
    \bottomrule
  \end{tabular}
  \smallskip

  La worklist est vide après le pas 13. Au pas 11, l'arête \og vrai\fg{}
  raffine $[0,5]$ en $[0,4]$ (\code{i} est entier, \code{i < 5} donne
  \code{i <= 4}) et l'arête \og faux\fg{} en $[5,5]$. Le slot \code{total} vaut
  $[0,5]$ : la valeur initiale $0$ et la valeur écrite $5$, réunies par join.
  C'est exactement la valeur qu'assère le test.
\end{exemple}

Les pas 4, 7 et 10 visitent $B_3$ pour rien : son entrée est vide à chaque
fois, mais elle passe de $[5,0]$ à $[5,1]$ puis $[5,2]$, et ces trois
représentations de $\Bot$ sont différentes pour \code{PartialEq}
(\cref{sec:interpretation-abstraite-concretisation}, piège sur les $\Bot$
multiples). C'est un coût, pas une faute.

\subsection{Le point fixe des tours de React}

Le même schéma existe à l'échelle du composant. La boucle externe de
\code{analyze_component_impl} (\cref{chap:point-fixe-composant}, et
\cref{chap:react-modele} pour sa lecture comme sémantique de React) exécute un
tour complet (rendu, mémos, effets, handlers) à partir du store courant, et
compare :

\rustsnippet[label={lst:interpretation-abstraite-convergence}]{src/engine/fixpoint.rs}{485}{497}{le test de post-point fixe de la boucle externe}

\code{new_state} est $F^\sharp(\code{state})$ : ce qu'un tour de React
abstrait produit à partir de \code{state}. Le test
\code{new_state.leq(&state)} est littéralement \og $F^\sharp(x) \lleq x$\fg{},
le test de post-point fixe du \cref{thm:post-point-fixe}. Tant qu'il échoue,
\code{state} est remplacé par \code{new_state} (les deux premières fois) puis
par un widening (\cref{sec:interpretation-abstraite-widening}).

\begin{piege}
  Le théorème ne garantit rien pour ce que le test \emph{ne regarde pas}. Le
  test de la boucle externe porte sur le store d'état et lui seul ; le tas et
  les mémos sont recalculés à chaque tour mais n'entrent pas dans la
  comparaison (\cref{chap:point-fixe-composant}). De même, au-delà de
  100 tours, le moteur fait un widening sans seuils puis sort \emph{sans}
  refaire le test (\src{src/engine/fixpoint.rs}{499}{512}) : l'état retenu
  n'est alors pas vérifié comme post-point fixe. L'argument de terminaison
  (\cref{sec:interpretation-abstraite-widening}) dit que ce garde-fou ne
  devrait jamais servir. Au commit \snapshotCommit{}, aucun test du dépôt ne
  le déclenche : c'est un filet de sécurité non exercé, pas un chemin validé.
\end{piege}

\section{Hauteur infinie et widening}
\label{sec:interpretation-abstraite-widening}

\subsection{Le problème}

Le \cref{thm:kleene} promet un nombre fini de pas si le treillis est de
hauteur finie. Les intervalles ne le sont pas
(\cref{fig:interpretation-abstraite-hasse-intervalles}), et
\cref{tab:interpretation-abstraite-tours} montre la conséquence : sur
\code{Runaway}, la suite $[0,1], [0,2], [0,3], \dots$ est une chaîne
strictement croissante qui ne s'arrête jamais. Le plus petit point fixe
abstrait existe, c'est $[0,+\infty]$, mais Kleene ne l'atteint qu'à la
limite. Il faut \emph{extrapoler} : deviner une borne qui majore toute la
suite, et vérifier ensuite que c'est un post-point fixe.

\begin{definition}{Widening, widening à seuils}{widening}
  Un \terme{widening} sur $L$ est un opérateur binaire $\widen$ tel que :
  \begin{enumerate}
    \item il majore le join : $a \join b \lleq a \widen b$ pour tous $a, b$ ;
    \item pour toute suite $y_0, y_1, \dots$, la suite
      $x_0 = y_0$, $x_{n+1} = x_n \widen y_{n+1}$ est ultimement stationnaire.
  \end{enumerate}
  Remplacer le join de Kleene par $\widen$ donne une suite qui se stabilise en
  un nombre fini de pas sur un post-point fixe, donc sur une réponse sound
  (\cref{thm:post-point-fixe}).

  Le widening des intervalles saute à l'infini dès qu'une borne grandit :
  \[ [a,b] \widen [c,d] = \big[\; c < a \;?\; -\infty : a,\;\; d > b \;?\; +\infty : b \;\big]. \]
  Le \terme{widening à seuils} $\widen_T$, pour un ensemble fini $T$ de
  nombres, saute au seuil le plus proche qui englobe la nouvelle borne, et à
  l'infini seulement s'il n'y en a pas :
  \[ [a,b] \widen_T [c,d] = [\ell, h] \quad\text{avec}\quad
     \begin{aligned}
       \ell &= \max\{t \in T \cup \{-\infty\} \mid t \le c\} \text{ si } c < a, \text{ sinon } a,\\
       h &= \min\{t \in T \cup \{+\infty\} \mid t \ge d\} \text{ si } d > b, \text{ sinon } b.
     \end{aligned} \]
\end{definition}

La première propriété garantit la correction : on ne perd jamais une valeur
que le join aurait gardée. La seconde garantit la terminaison : pour
$\widen$, chaque borne ne peut sauter qu'une fois (vers l'infini) ; pour
$\widen_T$, chaque borne ne peut prendre que des valeurs de $T \cup
\{\pm\infty\}$ après son premier saut, et ne fait que croître, donc au plus
$|T| + 1$ sauts par borne.

\begin{exemple}{Widening à la main}{widening-main}
  Avec $T = \{0, 1, 10\}$ :
  $[0,2] \widen [0,3] = [0, +\infty]$ (la borne haute grandit) ;
  $[0,2] \widen_T [0,3] = [0, 10]$ (le plus petit seuil $\ge 3$ est $10$) ;
  $[0,10] \widen_T [0,10] = [0,10]$ (rien ne grandit) ;
  $[0,10] \widen_T [0,11] = [0, +\infty]$ (aucun seuil $\ge 11$).
  Avec $T = \{0, 1\}$, $[0,2] \widen_T [0,3] = [0,+\infty]$.
\end{exemple}

\subsection{Ce que fait le code : \code{widen} et \code{widen_to}}

Le widening simple des intervalles est la formule ci-dessus, précédée des cas
$\Bot$ :

\rustsnippet[label={lst:interpretation-abstraite-widen}]{src/domains/impls/interval.rs}{84}{106}{le widening des intervalles}

Le widening à seuils (\adr{14}) remplace $\pm\infty$ par le seuil le plus
serré :

\rustsnippet[label={lst:interpretation-abstraite-widen-to}]{src/domains/impls/interval.rs}{108}{146}{le widening à seuils}

\begin{itemize}
  \item Une borne qui n'a pas grandi est conservée telle quelle
    (\code{self.lo}, \code{self.hi}).
  \item Une borne haute qui grandit prend le plus petit seuil
    $\ge$ \code{other.hi} : \code{filter} garde les candidats, \code{fold}
    avec \code{f64::min} à partir de $+\infty$ prend le minimum, et rend
    $+\infty$ si aucun seuil ne convient. La borne basse est symétrique. Les
    seuils n'ont pas besoin d'être triés.
  \item \code{is_int} reste vrai si les deux opérandes sont entiers ; une borne
    infinie ne rend pas l'intervalle non entier.
\end{itemize}

Aux autres composantes du produit \code{StateValue}, qui sont de hauteur
finie (treillis plats, ensembles bornés de chaînes, \code{Stability}), le
widening par défaut délègue au join
(\cref{lst:interpretation-abstraite-abstractdomain}, et
\src{src/domains/impls/mod.rs}{48}{50} pour les treillis plats).

\subsection{D'où viennent les seuils}

Les seuils sont récoltés une fois par composant, avant le point fixe
(\srcline{src/engine/fixpoint.rs}{280}) :

\rustsnippet[label={lst:interpretation-abstraite-thresholds}]{src/engine/fixpoint.rs}{1185}{1206}{la récolte des seuils}

Tous les littéraux numériques du rendu, des corps d'effets et de handlers, et
des initialiseurs de \code{useState} sont pris. L'ADR ne prévoyait que les
littéraux des gardes et des initialiseurs ; en prendre plus est sans danger :
l'ensemble reste fini (terminaison) et un seuil de plus n'est qu'une borne
candidate de plus (précision, jamais soundness). Pour \code{Guarded},
$T = \{0, 1, 10\}$ ; pour \code{Runaway}, $T = \{0, 1\}$.

\subsection{Retour aux compteurs}

Où le widening intervient-il ? Aux deux endroits où une suite peut croître
sans fin : les arcs arrière d'un CFG (\cref{lst:interpretation-abstraite-acfg-succ})
et la boucle des tours de React :

\rustsnippet[label={lst:interpretation-abstraite-outer-widen}]{src/engine/fixpoint.rs}{514}{529}{le widening de la boucle externe}

Le compteur \code{iteration} n'augmente que lorsqu'un tour a changé l'état. Les
deux premiers changements sont pris tels quels ; à partir du troisième
(\code{widen_threshold = 3}), l'état est élargi. Rejouons
\cref{tab:interpretation-abstraite-tours} avec le widening :

\begin{itemize}
  \item \code{Runaway}. Tours 0 et 1 : le slot passe à $[0,1]$ puis $[0,2]$.
    Tour 2 : le tour produit $[0,3]$, \code{iteration} vaut 3, donc
    $\code{state} = [0,2] \widen_{\{0,1\}} [0,3] = [0,+\infty]$. Tour 3 :
    l'effet écrit $[1,+\infty]$, le join donne $[0,+\infty]$, le test de
    post-point fixe réussit. Le slot convergé vaut $[0,+\infty]$, ce
    qu'épingle \code{unbounded_counter_is_flagged_and_grows_to_infinity}
    (\src{tests/widening_e2e.rs}{79}{92}).
  \item \code{Guarded}. Même début ; au tour 2,
    $[0,2] \widen_{\{0,1,10\}} [0,3] = [0,10]$. Tour 3 : la garde raffine
    $[0,10]$ en $[0,9]$, l'effet écrit $[1,10]$, le join rend $[0,10]$ : post-point
    fixe. Le test \code{guarded_counter_converges_bounded}
    (\src{tests/widening_e2e.rs}{94}{122}) épingle $[0,10]$.
\end{itemize}

La sortie de l'analyseur porte la trace de ces deux calculs. Avec
\texttt{-{}-info -{}-trace} (sortie élaguée des lignes \code{verified}) :

\begin{sortie}
$ reactant check /tmp/ch03/counter.tsx --info --trace
  Guarded  (2 hooks)  /tmp/ch03/counter.tsx
    info   widening-info  (line 13:2)  state `count` kept changing during analysis and was approximated to converge, so findings that depend on it may be imprecise
       → state `count` is written here [hook:1] (line 13:2)
       → the abstract value of state `count` kept growing and was widened at iteration 3
  Runaway  (2 hooks)  /tmp/ch03/counter.tsx
    warn   infinite-loop  [hook:0]  (line 5:2)  this effect keeps pushing state `count` (its deps do not provably gate it, so the effect can re-run every render) to new values on every run. Potential infinite render loop
       → state `count` is written here [hook:1] (line 5:2)
       → the abstract value of state `count` kept growing and was widened at iteration 3
    info   widening-info  (line 5:2)  state `count` kept changing during analysis and was approximated to converge, so findings that depend on it may be imprecise
       → state `count` is written here [hook:1] (line 5:2)
       → the abstract value of state `count` kept growing and was widened at iteration 3

⚠  1 warning(s) across 1 file(s).
\end{sortie}

Les deux composants ont subi un widening au tour où \code{iteration} a atteint 3
(l'Info \regle{widening-info} le dit, car un widening est une perte de
précision possible). Seul \code{Runaway} déclenche \regle{infinite-loop} :
son slot est non borné, celui de \code{Guarded} ne l'est pas. Le
\cref{chap:infinite-loop} détaille comment la règle lit cette différence ; le
niveau \Warning{} (et non \Error) s'expliquera à la
\cref{sec:interpretation-abstraite-must-may}.

\subsection{Le narrowing, et pourquoi \reactant{} s'en passe}

La théorie classique complète le widening par une phase descendante, le
\emph{narrowing} : une fois un post-point fixe $x$ atteint, on réapplique
$F^\sharp$ quelques fois, $x \meet F^\sharp(x)$, pour redescendre vers le plus
petit point fixe et récupérer la précision perdue par le saut à l'infini.

\begin{decision}[ADR-14 --- seuils oui, narrowing non]
  \adr{14} prévoyait les deux : widening à seuils (partie 1) et narrowing
  (partie 2). La partie 1 est implémentée, sur la boucle externe et sur les
  arcs arrière d'\code{analyze_cfg}. La partie 2 a été abandonnée
  (révision du 28 juin 2026) sur un constat empirique : le compteur gardé
  est déjà précis ($[0,10]$) grâce au raffinement de garde pendant la montée,
  et la boucle locale l'est devenue dès que le widening à seuils a été porté
  sur l'arc arrière interne ($[0,5]$). Les deux techniques ne récupèrent que
  des bornes \emph{littérales}, que les seuils capturent déjà en montant.
  Surtout, une phase descendante ne pourrait pas faire redescendre le store
  d'état, puisque les écritures s'y accumulent par join
  (\cref{sec:interpretation-abstraite-mises-a-jour}). Alternative écartée : un
  opérateur \code{narrow} dans \code{AbstractDomain} et une boucle
  descendante bornée ; il reste une infrastructure différée, à reconsidérer
  seulement si des bornes symboliques deviennent nécessaires.
\end{decision}

Ce que ni les seuils ni le narrowing ne savent faire, c'est une borne
symbolique. Remplaçons le littéral par une prop :

\begin{lstlisting}[language=TypeScript,caption={Une garde à borne symbolique (\file{/tmp/ch03/limit.tsx}, élagué de sa ligne d'import)},label={lst:interpretation-abstraite-limit}]
export function Limited({ limit }: { limit: number }) {
  const [count, setCount] = useState(0);
  useEffect(() => {
    if (count < limit) setCount(count + 1);
  }, [count, limit]);
  return <div>{count}</div>;
}
\end{lstlisting}

\begin{sortie}
$ reactant check /tmp/ch03/limit.tsx
  Limited  (2 hooks)  /tmp/ch03/limit.tsx
    warn   infinite-loop  [hook:0]  (line 5:2)  this effect keeps pushing state `count` (its deps do not provably gate it, so the effect can re-run every render) to new values on every run. Potential infinite render loop
\end{sortie}

Pour une valeur finie de \code{limit}, ce composant s'arrête. Mais la garde
\code{count < limit} compare deux variables : le raffinement de garde ne
reconnaît que \og variable comparée à un littéral\fg{}
(\cref{chap:cfg-analyzer-dominance}), aucun seuil ne vaut \code{limit}, et le
widening saute à $+\infty$. Le \Warning{} est un faux positif, toléré par
l'invariant. Le corriger demanderait un domaine \emph{relationnel} (qui
retient $\code{count} \le \code{limit}$), ce que \reactant{} n'a pas
(\cref{sec:interpretation-abstraite-reactant}).

\begin{piege}
  Le widening n'est pas monotone et dépend de l'ordre des calculs :
  $[0,1] \lleq [0,2]$, et pourtant $[0,1] \widen [0,2] = [0,+\infty]$ alors
  que $[0,2] \widen [0,2] = [0,2]$ ; un argument gauche plus précis donne un
  résultat moins précis. Élargir au tour 2 ou au tour 3 ne donne donc pas
  forcément le même résultat. C'est
  sans conséquence sur la soundness, qui ne repose que sur le test de
  post-point fixe final (\cref{thm:post-point-fixe}). Mais un changement
  anodin (un seuil de plus, un \code{widen_threshold} différent, un ordre de
  visite des blocs différent) peut changer la précision, donc faire
  apparaître ou disparaître un faux positif.
\end{piege}

\section{Must, may, et la polarité d'un fait}
\label{sec:interpretation-abstraite-must-may}

\subsection{Deux setters dans le rendu}

Appeler un setter pendant le rendu est un défaut (\cref{chap:react-modele}).
Comparons deux façons de le faire :

\begin{lstlisting}[language=TypeScript,caption={Un setter dans le rendu, inconditionnel puis conditionnel (\file{/tmp/ch03/setter.tsx}, élagué de sa ligne d'import)},label={lst:interpretation-abstraite-setter}]
export function Unconditional() {
  const [n, setN] = useState(0);
  setN(n + 1);
  return <p>{n}</p>;
}

export function Conditional({ flag }: { flag: boolean }) {
  const [n, setN] = useState(0);
  if (flag) setN(1);
  return <p>{n}</p>;
}
\end{lstlisting}

\begin{sortie}
$ reactant check /tmp/ch03/setter.tsx
  Conditional  (1 hooks)  /tmp/ch03/setter.tsx
    warn   setter-in-render  [hook:0]  (line 11:12)  setter `setN` called directly in the render body, move this call into a useEffect or an event handler
       (1 trace step(s), rerun with --trace)
  Unconditional  (1 hooks)  /tmp/ch03/setter.tsx
    error  setter-in-render  [hook:0]  (line 5:2)  setter `setN` called directly in the render body, move this call into a useEffect or an event handler
       (1 trace step(s), rerun with --trace)

⚠  1 error(s), 1 warning(s) across 1 file(s).
\end{sortie}

Même règle, deux niveaux. Dans \code{Unconditional}, \emph{tout} rendu
appelle \code{setN} : le défaut a lieu dès que le composant s'affiche. Dans
\code{Conditional}, il n'a lieu que si \code{flag} est vrai, ce que
l'analyseur ne sait pas (\code{flag} est une prop, donc $\Top$). Le premier
fait est vrai sur \emph{tous} les chemins, le second sur \emph{certains}.

\subsection{Pourquoi une sur-approximation ne suffit pas}

Tout ce qui précède calcule des sur-approximations : l'ensemble abstrait
contient tous les comportements possibles, et peut-être d'autres. Une
sur-approximation prouve des \emph{absences} (\og \code{count} ne dépasse
jamais 10\fg{}), jamais des \emph{présences}. Si l'abstrait dit
\og \code{setN} est peut-être appelé\fg{}, c'est peut-être à cause d'un chemin
impossible. Pour affirmer qu'un défaut \emph{a lieu}, il faut l'autre sens :
un fait vrai de toutes les exécutions, donc une sous-approximation de ce qui
arrive à coup sûr.

\begin{definition}{Fait must, fait may, polarité}{must-may}
  Un fait est \terme{must} s'il est établi pour \emph{toutes} les exécutions
  concrètes que l'abstraction décrit (sur tous les chemins, à tous les tours).
  Il est \terme{may} s'il vaut pour \emph{au moins une} exécution que
  l'abstraction admet, qui peut être une exécution impossible introduite par
  l'approximation. $\Top$ (l'inconnu) n'est ni l'un ni l'autre ; lu comme
  réponse à la question \og le défaut a-t-il lieu ?\fg{}, il ne peut compter
  que comme may. La \terme{polarité} d'un fait, d'une relation ou d'une
  colonne de relation (\cref{def:relation}) dit lequel des deux elle porte :
  exacte, must, may ou $\Top$.
\end{definition}

Les deux polarités se calculent avec les mêmes outils, en sens inverse. Une
analyse may joint aux jonctions (union des possibilités) et part de $\Bot$ :
plus petit point fixe. Une analyse must intersecte aux jonctions (ce qui est
vrai de \emph{chaque} prédécesseur) et part de \og tout est vrai\fg{} : plus
grand point fixe. Nous en verrons un exemple avec la dominance
(\cref{sec:interpretation-abstraite-dominance}).

\subsection{Ce que fait le code : des verdicts typés par leur polarité}

La couche règles (\cref{chap:api-regles}) inscrit la polarité dans les types
(\adr{21}). Une primitive must rend un \code{MustResult} :

\rustsnippet[label={lst:interpretation-abstraite-mustresult}]{src/rules/api/query.rs}{111}{135}{les verdicts polarisés}

\begin{itemize}
  \item \code{All} : établi sur tous les chemins. C'est le seul cas qui
    transporte un \code{Certified} (\cref{def:certified}), le jeton sans lequel
    aucune règle ne peut émettre une \Error.
  \item \code{Some} : établi sur certains chemins seulement, un fait may, sans
    jeton.
  \item \code{None} : aucune évidence.
  \item \code{May<T>} enveloppe un fait may ; aucune fonction ne le
    transforme en \code{Certified}.
\end{itemize}

Le contrat de toutes les primitives tient en trois clauses
(\src{src/rules/api/query.rs}{417}{421}) : chaque primitive rend un verdict
typé par sa polarité ; $\Top$ se replie en may à l'intérieur ; seules les
primitives must frappent un \code{Certified}. La deuxième clause est celle
qui protège la soundness des \Error : un inconnu ne devient jamais une
certitude.

Voici comment \regle{setter-in-render} fabrique la différence observée plus
haut :

\rustsnippet[label={lst:interpretation-abstraite-sir-proof}]{src/rules/impls/setter_in_render.rs}{147}{162}{\regle{setter-in-render} : l'\Error{} demande deux musts}

La preuve a deux conjoints must : l'appel a lieu pendant le rendu, et il
écrit à coup sûr (\code{certain_write}) ; puis son bloc domine toutes les
sorties du rendu (\code{exit_dom.certify}, \cref{sec:interpretation-abstraite-dominance}).
Si l'un manque, \code{proof} vaut \code{None} et le diagnostic reste un
\Warning. Dans \code{Conditional}, le bloc \code{then} ne domine pas le
\code{return} : pas de jeton, \Warning.

\begin{piege}
  Un fait must sur \emph{une partie} de la conclusion ne suffit pas. La
  doctrine du projet (\file{CLAUDE.md}) exige, pour une \Error, la preuve
  \og de \emph{toute} la conclusion, pas d'un seul conjoint\fg{}.
  \regle{infinite-loop} en donne un exemple : pour \code{Runaway}, la
  croissance du slot jusqu'à $+\infty$ est le résultat d'un widening,
  c'est-à-dire d'une extrapolation sur-approximante ; elle ne prouve pas que
  l'effet re-tourne indéfiniment. La règle n'émet donc qu'un \Warning{} ;
  seule la preuve d'un cycle d'effets (\cref{chap:infinite-loop}) atteint
  l'\Error. L'écart entre divergence de valeur et divergence de référence est
  suivi par l'\issue{144}.
\end{piege}

\section{Soundness}
\label{sec:interpretation-abstraite-soundness}

Nous pouvons maintenant donner à la promesse de l'introduction
(\cref{def:introduction-soundness}) sa forme exacte.

\begin{definition}{Soundness}{soundness}
  Une analyse est \terme{sound} lorsque ce qu'elle calcule
  \emph{sur-approxime} la sémantique concrète : pour tout programme,
  $\lfp F \subseteq \gam(a)$ où $a$ est le résultat abstrait. Pour un outil
  de diagnostic, cela se traduit ainsi : si le programme peut exhiber un
  défaut qu'une règle vise, l'analyse le signale (à un niveau ou à un autre),
  ou signale explicitement qu'elle n'a pas pu regarder. Un défaut réel non
  signalé, un \terme{faux négatif} (FN), est interdit ; un signalement sans
  défaut réel, un \terme{faux positif} (FP), est toléré.
\end{definition}

Chaque niveau de diagnostic en tire une contrainte différente.

\begin{itemize}
  \item \textbf{Le silence.} Une règle muette affirme l'absence du défaut. Cette
    affirmation est une propriété de sur-approximation : elle n'est vraie que
    si tous les calculs amont ont été corrects (\cref{def:concretisation}) et
    ont atteint un post-point fixe (\cref{thm:post-point-fixe}). Sous
    \texttt{-{}-info}, les lignes \code{verified} rendent ce silence visible.
  \item \textbf{\Error.} Un défaut certain dès que le code s'exécute. C'est une
    affirmation must, qui ne peut pas venir d'une sur-approximation : elle
    demande un \code{Certified} (\cref{sec:interpretation-abstraite-must-may}).
  \item \textbf{\Warning.} Un défaut possible (fait may), ou un fait certain
    au coût incertain. Tout ce que la sur-approximation ne permet pas
    d'exclure et que rien ne permet de certifier tombe ici.
  \item \textbf{\Info.} Une limite de l'analyse, ou un motif qui a l'air
    voulu. Quand l'analyse n'a pas pu lire du code (un inlining tronqué, par
    exemple), elle le dit par une \Info{} et suspend ses assurances
    \code{verified} (\cref{chap:regles-effets}) : un silence qui ne vaut plus
    preuve est retiré plutôt que maintenu.
\end{itemize}

\begin{exemple}{Un faux positif toléré}{fp-dead}
  \begin{lstlisting}[language=TypeScript]
export function Dead() {
  const [n, setN] = useState(0);
  const k = 3;
  useEffect(() => {
    if (k > 10) setN(n + 1);
  }, [n]);
  return <div>{n}</div>;
}
  \end{lstlisting}
  \begin{sortie}
$ reactant check /tmp/ch03/dead.tsx
  Dead  (2 hooks)  /tmp/ch03/dead.tsx
    warn   infinite-loop  [hook:0]  (line 6:2)  this effect keeps pushing state `n` (its deps do not provably gate it, so the effect can re-run every render) to new values on every run. Potential infinite render loop
  \end{sortie}
  La branche est morte : \code{k} vaut 3. Le raffinement de \code{k > 10}
  donne bien $\Bot$ à \code{k} sur la branche vraie, mais un raffinement vers
  $\Bot$ ne rend pas le bloc inatteignable : le bloc est exécuté, seule la
  variable gardée y vaut $\Bot$, et \code{setN(n + 1)} écrit une vraie valeur
  (\cref{chap:cfg-analyzer-dominance}). Le \Warning{} est un FP : imprécis,
  sound.
\end{exemple}

\begin{remarque}
  La soundness de \reactant{} n'est pas un théorème prouvé une fois pour
  toutes ; c'est un invariant tenu opération par opération : join et widening
  majorent le join ($\widen_T$ \og $\sqsupseteq$ hull\fg{}), une variable
  inconnue vaut $\Top$, une écriture de slot est jointe et jamais remplacée, un
  appel opaque rend $\Top$, un must ne compte pas un chemin d'exception comme
  couvert. Le \cref{chap:point-fixe-composant} en dresse l'inventaire. Là où
  la soundness dépend d'une hypothèse du modèle, une règle dédiée surveille
  l'hypothèse (soundness en couches, \cref{def:soundness-couches}). Les
  décisions de ne \emph{pas} changer certains comportements pour des raisons
  de soundness sont consignées dans \adr{20}.
\end{remarque}

\begin{piege}
  \og Tolérer les FP\fg{} ne veut pas dire \og les ignorer\fg{}. Un FP au
  niveau \Error{} est une faute : l'\Error{} promet un défaut certain, et le
  contrat de la \cref{sec:interpretation-abstraite-must-may} existe pour
  l'empêcher. Les FP tolérés sont des \Warning{} (ou des \Info) produits par
  l'imprécision de la sur-approximation.
\end{piege}

\section{CFG, dominance, \og sur tous les chemins\fg{}}
\label{sec:interpretation-abstraite-dominance}

\subsection{Un hook après un \code{return} anticipé}

\begin{lstlisting}[language=TypeScript,caption={Un hook sauté par un \code{return} anticipé (\file{/tmp/ch03/hook.tsx}, élagué de sa ligne d'import)},label={lst:interpretation-abstraite-hook}]
export function Profile({ user }: { user: string | null }) {
  if (user === null) {
    return <p>anonymous</p>;
  }
  const [name, setName] = useState(user);
  useEffect(() => {
    document.title = name;
  }, [name]);
  return <p>{name}</p>;
}
\end{lstlisting}

\begin{sortie}
$ reactant check /tmp/ch03/hook.tsx --trace
  Profile  (2 hooks)  /tmp/ch03/hook.tsx
    error  conditional-hook  [hook:0]  (line 7:8)  this hook is called conditionally (not on every render path)
       → guarded by a condition evaluated here, so some render paths skip the hook (line 4:6)
    error  conditional-hook  [hook:1]  (line 8:2)  this hook is called conditionally (not on every render path)
       → guarded by a condition evaluated here, so some render paths skip the hook (line 4:6)

⚠  2 error(s) across 1 file(s).
\end{sortie}

Les règles des hooks exigent qu'un hook soit appelé à chaque rendu. Sur le CFG
du rendu (\cref{fig:interpretation-abstraite-cfg-hook}), cela veut dire :
\emph{tout} chemin de l'entrée à une sortie passe par le bloc du hook. Ce
\og tout\fg{} est une quantification sur une infinité potentielle de chemins ;
la dominance la ramène à un calcul fini.

\begin{figure}[htbp]\centering
\begin{tikzpicture}[node distance=7mm and 12mm]
  \node[cfgbloc] (b0) {$B_0$ (entrée)\\branch user === null};
  \node[cfgbloc,below left=of b0] (b1) {$B_1$\\return <p>anonymous</p>};
  \node[cfgbloc,below right=of b0] (b2) {$B_2$\\useState(user)\\useEffect(…)\\return <p>{name}</p>};
  \draw[fleche] (b0) -- node[etiquette,left]{vrai} (b1);
  \draw[fleche] (b0) -- node[etiquette,right]{faux} (b2);
\end{tikzpicture}
\caption{CFG schématique du rendu de \code{Profile}. $\mathrm{Dom}(B_1) = \{B_0,
B_1\}$ et $\mathrm{Dom}(B_2) = \{B_0, B_2\}$ : $B_2$, qui porte les hooks, ne
domine pas la sortie $B_1$.}
\label{fig:interpretation-abstraite-cfg-hook}
\end{figure}

\begin{definition}{Dominance, post-dominance}{dominance}
  Dans un CFG d'entrée $e$, un bloc $a$ \terme{domine} un bloc $b$ si tout
  chemin de $e$ à $b$ passe par $a$. On note $\mathrm{Dom}(b)$ l'ensemble des
  dominateurs de $b$ ; c'est la plus grande solution du système
  \[ \mathrm{Dom}(e) = \{e\}, \qquad
     \mathrm{Dom}(b) = \{b\} \cup \bigcap_{p \in \pred(b)} \mathrm{Dom}(p). \]
  Un ensemble de blocs $S$ est \terme{sur tous les chemins} (il post-domine
  l'entrée) si tout chemin de $e$ à une sortie rencontre $S$. Un bloc
  s'exécute à chaque passage dans le CFG exactement quand il domine toutes
  les sorties atteignables.
\end{definition}

Le système de la dominance est une analyse must : l'opérateur de jonction est
l'intersection, et l'on part de \og tous les blocs dominent tout\fg{} pour
descendre. C'est donc le \emph{plus grand} point fixe qu'on cherche, par une
itération duale de Kleene (\cref{def:point-fixe}) : le treillis est celui des
parties de blocs ordonné par $\supseteq$, fini, et l'itération termine.

\subsection{Ce que fait le code}

\code{compute_dominators} est cet algorithme itératif classique, parcouru en
ordre postfixe inverse (RPO) :

\rustsnippet[label={lst:interpretation-abstraite-dominators}]{src/engine/dominance.rs}{22}{57}{le plus grand point fixe de la dominance}

Chaque bloc autre que l'entrée part de \code{all_blocks} (lignes 16 à 20,
hors extrait) et descend par intersection. Un bloc inatteignable n'est pas
dans le RPO et garde \og tous les blocs\fg{} : il est dominé par tout, ce qui
est neutre pour un fait must (il ne s'exécute jamais). Pour éviter de refaire
ce calcul à chaque question, \code{DominatorTree}
(\src{src/engine/dominance.rs}{94}{111}) le mémorise.

La couche règles consomme la dominance sous une forme typée, qui choisit les
sorties une fois pour toutes :

\rustsnippet[label={lst:interpretation-abstraite-exitdom}]{src/rules/api/query.rs}{672}{706}{\code{ExitDominance} : négation may et certification must}

\begin{itemize}
  \item \code{certify} rend \code{All} (avec un \code{Certified}) si le bloc
    domine \emph{toutes} les sorties : c'est l'arme de \regle{setter-in-render}
    (\cref{lst:interpretation-abstraite-sir-proof}).
  \item \code{may_be_skipped} en est la négation : il existe une sortie que le
    bloc ne domine pas. Pour \code{Profile}, $B_2$ ne domine pas $B_1$ ; le hook
    est sauté sur au moins un chemin. Le défaut visé est structurel (la règle
    des hooks porte sur la forme du code, pas sur les valeurs) : l'existence
    de ce chemin dans le CFG suffit à l'établir, et
    \regle{conditional-hook} l'émet en \Error{} (\cref{def:hook-conditionnel}).
  \item Un CFG sans sortie atteignable ne prouve rien dans aucun sens :
    \code{certify} rend \code{None}, et \code{may_be_skipped} rend faux.
\end{itemize}

\begin{piege}
  Les sorties comptées sont les \code{return} \emph{atteignables}
  (\src{src/rules/api/query.rs}{652}{671}). Le lowering scelle la fin d'un
  corps par un \code{return undefined} (\adr{25}) ; quand les deux branches
  d'un \code{if}/\code{else} ont déjà retourné, cette queue est orpheline. La
  compter ferait croire que tout hook antérieur est sauté, au niveau
  \Error{} : un FP interdit par la \cref{sec:interpretation-abstraite-soundness}.
\end{piege}

\code{on_all_paths} (\src{src/engine/dominance.rs}{69}{92}) répond à la
question voisine \og ces blocs sont-ils sur tous les chemins ?\fg{} par un
parcours qui évite l'ensemble : atteindre une feuille (un \code{return}, mais
aussi un \code{throw}) sans le rencontrer, c'est exhiber un chemin qui
s'échappe. Sa polarité est must : un chemin d'exception compte comme un
chemin qui évite l'ensemble. Une troisième implémentation du même
quantificateur existe, l'analyse must avant de \code{must_setter_on_all_paths}
(\src{src/rules/api/query.rs}{576}{621}), qui calcule elle aussi un plus grand
point fixe booléen ; le \cref{chap:api-regles} la présente.

\section{Intra- et inter-procédural, inlining}
\label{sec:interpretation-abstraite-inlining}

Tout ce qui précède opère sur un CFG. Or un composant appelle des fonctions :
des hooks personnalisés, des utilitaires, des composants enfants. Une analyse
\terme{intra-procédurale} traite chaque fonction isolément et remplace un
appel par ce qu'elle sait sans le lire, c'est-à-dire $\Top$ : sound, mais
aveugle. Une analyse \terme{inter-procédurale} suit les appels. Deux familles
de techniques existent : calculer un \emph{résumé} de chaque fonction,
réutilisé à chaque appel ; ou \terme{inliner}, c'est-à-dire greffer le corps
de la fonction appelée au site d'appel et analyser le tout comme un seul CFG.
L'inlining est la plus précise (chaque appel est analysé dans son contexte)
et la plus coûteuse (le CFG grossit, la récursion doit être coupée).

\begin{lstlisting}[language=TypeScript,caption={Un hook personnalisé qui boucle (\file{/tmp/ch03/custom.tsx}, élagué de sa ligne d'import)},label={lst:interpretation-abstraite-custom}]
function useTicker(start: number) {
  const [t, setT] = useState(start);
  useEffect(() => {
    setT(t + 1);
  }, [t]);
  return t;
}

export function Clock() {
  const t = useTicker(0);
  return <p>{t}</p>;
}
\end{lstlisting}

\begin{sortie}
$ reactant check /tmp/ch03/custom.tsx
  Clock  (1 hooks)  /tmp/ch03/custom.tsx
    warn   infinite-loop  [hook:1]  (line 5:2)  this effect keeps pushing state `t` (its deps do not provably gate it, so the effect can re-run every render) to new values on every run. Potential infinite render loop
       (2 trace step(s), rerun with --trace)
\end{sortie}

\code{Clock} ne contient qu'un hook, l'appel à \code{useTicker}. Le
diagnostic désigne pourtant \code{[hook:1]}, à la ligne 5 de
\code{useTicker} : le moteur a greffé le corps du hook dans le rendu de
\code{Clock}, décalé ses labels de hooks, et analysé le compteur comme s'il
était écrit dans le composant. \reactant{} fait ce choix partout où il
dispose du source.

\begin{itemize}
  \item \textbf{Hooks personnalisés} : greffés entiers au site d'appel, labels
    décalés, paramètres substitués (\code{expand_custom_hooks},
    \src{src/engine/fixpoint.rs}{884}{1142}). Un hook qui s'appelle lui-même,
    ou un second appel du même hook dans le même composant, reste opaque
    ($\Top$). Le commentaire du code l'avoue (\og \emph{second call stays opaque
    FN, not FP}\fg{}) ; la sortie le compense : appeler \code{useTicker} deux
    fois (\file{/tmp/ch03/twice.tsx}) produit une \Info{}
    \regle{analysis-limit} sur le second appel (\og \emph{hook `useTicker` was
    not found in the registry}\fg{}, message trompeur puisque le hook y est) et
    la ligne \code{suspended analysis-limit 9 passing check(s) withheld}.
  \item \textbf{Hooks de bibliothèque sans source} : un registre de résumés
    (\code{SummaryRegistry}) donne leur valeur abstraite ; c'est la technique
    par résumé, réservée à ce cas.
  \item \textbf{Utilitaires} : greffés en position d'instruction, au plus une
    fois par nom et par CFG, dans un budget de \code{max_inline_depth} = 8
    greffes par CFG (\src{src/engine/fixpoint.rs}{1611}{1634}) ; un budget
    épuisé laisse les appels restants opaques et suspend les assurances du
    composant (\cref{sec:interpretation-abstraite-soundness}, rubrique \Info).
  \item \textbf{Composants enfants} : analysés de haut en bas avec les props
    que le parent leur passe, et leurs écritures dans l'état du parent
    remontent par un store partagé (\adr{12}).
\end{itemize}

Les mécanismes de greffe (\emph{splice}, renommage, identités) sont l'objet du
\cref{chap:splice-inlining}, l'analyse entre composants et entre fichiers
celui du \cref{chap:inter-composants}. Retenons ici leur statut théorique :
l'inlining ne change rien au cadre de ce chapitre, il agrandit le CFG sur
lequel on calcule le point fixe ; et chaque endroit où il renonce (récursion,
budget, callee inconnu) retombe sur $\Top$, donc reste sound.

\section{Mises à jour fortes et faibles}
\label{sec:interpretation-abstraite-mises-a-jour}

\begin{lstlisting}[language=TypeScript,caption={Deux écritures successives (\file{/tmp/ch03/weak.tsx}, élagué de sa ligne d'import)},label={lst:interpretation-abstraite-weak}]
export function Twice() {
  const [n, setN] = useState(0);
  useEffect(() => {
    let x = 1;
    x = 5;
    setN(x);
    setN(7);
  }, []);
  return <div>{n}</div>;
}
\end{lstlisting}

Après \code{x = 5}, l'ancienne valeur de \code{x} est perdue : il est exact
de \emph{remplacer} $[1,1]$ par $[5,5]$. C'est une \terme{mise à jour forte}.
Elle n'est correcte que si l'écriture désigne un unique lieu de stockage
concret, qu'elle a certainement lieu, et que l'abstraction ne décrit pas
plusieurs instants à la fois. Quand l'une de ces conditions manque, il faut
une \terme{mise à jour faible} : joindre la nouvelle valeur à l'ancienne,
$v \gets v \join v'$, ce qui dit \og ce lieu vaut l'ancienne valeur ou la
nouvelle\fg{}.

\reactant{} emploie les deux, selon ce que le lieu abstrait représente.
\begin{itemize}
  \item \textbf{Variables locales : forte.} L'environnement est calculé bloc
    par bloc ; une affectation remplace la liaison
    (\code{AbstractEnv::extend} est un \code{insert},
    \src{src/domains/stores/abstract_env.rs}{100}{103}). Les chemins sont
    réunis par join aux jonctions du CFG.
  \item \textbf{Slots d'état : faible.} Un slot abstrait décrit \emph{toutes}
    les valeurs que l'état prend au fil de \emph{tous} les tours de React,
    pas sa valeur à un instant. Écrire ne peut donc qu'ajouter :
\end{itemize}

\rustsnippet[label={lst:interpretation-abstraite-update}]{src/domains/stores/state_store.rs}{29}{33}{l'écriture d'un slot est un join}

\begin{itemize}
  \item \textbf{Champs du tas : faible.} Un objet abstrait est identifié par son
    site d'allocation (\cref{def:site}) et représente tous les objets créés à ce
    site ; une écriture \code{obj.f = v} peut n'avoir lieu que sur un chemin, et
    \code{obj} peut désigner plusieurs sites. Le code joint donc, et le dit :
    \og \emph{join, never replace}\fg{}
    (\src{src/domains/interp/interpreter.rs}{150}{154}).
\end{itemize}

Sur \code{Twice}, le slot \code{n} reçoit $[0,0]$ (amorçage), puis
$[5,5]$, puis $[7,7]$ : il vaut $[0,0] \join [5,5] \join [7,7] = [0,7]$
(calcul fait d'après \cref{lst:interpretation-abstraite-update}). En React,
les deux appels sont regroupés (batching) et le rendu suivant voit 7 ;
l'abstrait contient $0$ et $7$, et aussi $1, \dots, 6$ : une perte de
précision, pas d'information fausse. C'est aussi pourquoi le narrowing ne
pourrait pas faire redescendre un slot
(\cref{sec:interpretation-abstraite-widening}) : aucune écriture n'est jamais
retirée.

\begin{sortie}
$ reactant check /tmp/ch03/weak.tsx
  Twice  (2 hooks)  /tmp/ch03/weak.tsx
    warn   unnecessary-rerender  [hook:0]  (line 5:2)  mount-only effect sets state `n` to a constant different from its initial value, which costs one extra rerender on mount; consider initialising directly with the target value
       (1 trace step(s), rerun with --trace)
\end{sortie}

La règle qui répond ici ne lit pas le slot joint ; elle parcourt les appels de
setter de l'effet de montage et évalue leurs arguments un par un
(\src{src/rules/impls/unnecessary_rerender.rs}{110}{144}, et
\cref{chap:regles-etat}). La
mise à jour faible est la bonne abstraction du store, pas la seule source
d'information des règles : quand une question porte sur \emph{une} écriture,
le moteur la conserve à part, site par site (\cref{def:slot-writer}).

\begin{piege}
  Rendre fortes les écritures d'état \og parce que le dernier \code{setN}
  gagne\fg{} serait un faux négatif en puissance : le slot ne décrirait plus
  les valeurs vues par les rendus intermédiaires, ni celles écrites par un
  autre effet ou un handler au même tour. La règle \og dernier set gagne\fg{}
  vaut pour un tour concret donné ; le slot abstrait les résume tous.
\end{piege}

\section{Ce que \reactant{} fait, et ne fait pas}
\label{sec:interpretation-abstraite-reactant}

Un analyseur par interprétation abstraite se caractérise par une poignée de
choix : ce qu'il distingue (flot, chemins, contextes), ce que ses domaines
savent exprimer, et comment il force la terminaison.
\Cref{tab:interpretation-abstraite-choix} résume ceux de \reactant{} ;
chacun a été rencontré dans ce chapitre.

\begin{table}[htbp]\centering\small
\begin{tabularx}{\linewidth}{@{}lX@{}}
  \toprule
  Dimension & Choix de \reactant{} \\
  \midrule
  Sensibilité au flot (variables) & Oui : un environnement par entrée et sortie de bloc (\code{entry_envs}, \code{exit_envs}), mises à jour fortes. \\
  Sensibilité au flot (état) & Non : un seul store par passe (\code{state_out}), partagé par tous les blocs, et un seul store pour tous les tours ; mises à jour faibles. \\
  Sensibilité aux chemins & Non, sauf le raffinement de garde : les chemins sont joints aux jonctions ; une garde \og variable comparée à un littéral\fg{} raffine chaque branche, sans jamais rendre une branche morte. \\
  Forme de l'IR & Pas de SSA : les variables sont des noms, une affectation rebinde le nom dans l'environnement (\cref{chap:ir}). \\
  Domaines numériques & Non relationnels : un intervalle par variable ou par slot ; aucune relation entre deux variables (d'où le FP de \code{Limited}). \\
  Inter-procédural & Inlining des hooks, des utilitaires et des enfants, coupé à la récursion et par un budget ; résumés pour les hooks de bibliothèque ; $\Top$ ailleurs. \\
  Terminaison & Widening à seuils après 3 changements, sur les arcs arrière et entre tours ; widening sans seuils et arrêt au-delà de 100 tours. Pas de narrowing (\adr{14}). \\
  Polarité & Sur-approximation (may) pour le silence et les \Warning ; faits must certifiés (dominance, analyses must avant) pour les \Error. \\
  \bottomrule
\end{tabularx}
\caption{Les choix d'interprétation abstraite de \reactant{}.}
\label{tab:interpretation-abstraite-choix}
\end{table}

Deux de ces choix méritent une phrase de justification. L'état est
insensible au flot parce que la question que posent les règles porte sur
l'ensemble des valeurs qu'un slot prend au cours de la vie du composant, pas
sur sa valeur en un point du rendu : React ne change un slot qu'entre deux
rendus, et le résumé de tous les tours est précisément ce qu'il faut. Les
domaines sont non relationnels parce que les bornes que \adr{14} a
rencontrées sont des littéraux, que les seuils capturent déjà ; un domaine
relationnel (octogones, polyèdres) ne se justifierait que pour des bornes
symboliques comme celle de \code{Limited}, cas que l'ADR laisse ouvert.

Les chapitres suivants reprennent chaque ligne du tableau à la
\emph{taille réelle} du code : les domaines de base
(\cref{chap:treillis-simples}), le produit \code{StateValue} et la stabilité
(\cref{chap:statevalue-stabilite}), le transfert et les stores
(\cref{chap:transfert-stores}), la worklist et la dominance
(\cref{chap:cfg-analyzer-dominance}), la boucle externe
(\cref{chap:point-fixe-composant}).

\begin{resume}
\begin{itemize}
  \item Aucun analyseur ne peut calculer exactement les comportements d'un
    composant (taille, indécidabilité). \reactant{} \emph{sur-approxime} :
    l'abstrait contient tout le concret. Son silence est une preuve ; ses
    erreurs sont des faux positifs, jamais des faux négatifs
    (\cref{def:soundness}).
  \item Un domaine abstrait est un treillis (\cref{def:treillis}) : ordre
    partiel (\code{PartialOrd}), join, meet, $\Bot$, $\Top$. Sa
    concrétisation $\gam$ dit ce que signifie chaque valeur ; une opération
    est correcte si elle ne perd aucune valeur concrète
    (\cref{def:concretisation}). Trait \code{AbstractDomain}
    (\file{src/domains/mod.rs}), \code{BoolVal}, \code{flat_lattice!},
    \code{Interval}.
  \item Le résultat d'une analyse est un post-point fixe d'équations de flot
    (\cref{def:point-fixe}) ; tout post-point fixe d'une fonction abstraite
    correcte est sound (\cref{thm:post-point-fixe}). \code{analyze_cfg} est
    une worklist FIFO sur un CFG ; la boucle externe de
    \code{analyze_component_impl} teste \code{new_state.leq(&state)}.
  \item Les intervalles sont de hauteur infinie : le widening à seuils
    (\cref{def:widening}, \code{Interval::widen_to}, seuils récoltés par
    \code{collect_thresholds}) force la terminaison après
    \code{widen_threshold = 3} changements. Pas de narrowing (\adr{14}) : le
    raffinement de garde et les seuils récupèrent les bornes littérales, pas
    les bornes symboliques.
  \item Une \Error{} demande un fait \emph{must} (\cref{def:must-may}) ;
    \code{MustResult::All} est le seul porteur de \code{Certified} ; $\Top$ se
    replie en may. La dominance (\cref{def:dominance},
    \file{src/engine/dominance.rs}, \code{ExitDominance}) est l'analyse must
    de référence, un plus grand point fixe.
  \item L'environnement est sensible au flot et mis à jour fortement ; l'état
    et le tas sont insensibles au flot et mis à jour faiblement ; l'inlining
    rend l'analyse inter-procédurale, et chaque renoncement retombe sur
    $\Top$.
\end{itemize}
Le \cref{chap:tour-architecture} situe ces mécanismes dans l'architecture du
code ; la partie III (\cref{chap:treillis-simples} à
\cref{chap:transfert-stores}) reprend les domaines un par un, la partie IV
(\cref{chap:cfg-analyzer-dominance} et suivants) le moteur.
\end{resume}

\section*{Exercices}
\addcontentsline{toc}{section}{Exercices}

\begin{exercice}{Le treillis plat}{flat}
  (a) Avec \cref{lst:interpretation-abstraite-flat}, que rend
  \code{partial_cmp} sur $(\Bot, \Top)$, sur $(\Top, \Bot)$, sur
  $(\code{True}, \code{False})$ ? (b) Calculer
  $\code{One(c,0)} \join \code{One(c,1)}$ et $\code{One(c,0)} \meet \Top$ dans
  \code{SetterVal}. (c) Pourquoi le widening d'un treillis plat peut-il être
  le join ?
\end{exercice}
\begin{solution}
  (a) $(\Bot,\Top)$ : les éléments diffèrent, le deuxième bras
  s'applique (\og $\Bot$ à gauche\fg{}), \code{Some(Less)} ;
  $(\Top,\Bot)$ : ni égalité, ni deuxième bras, le troisième donne
  \code{Some(Greater)} ; $(\code{True},\code{False})$ : aucun bras,
  \code{None}. (b) Deux éléments intermédiaires distincts : $\Top$. Le meet
  avec $\Top$ rend l'autre opérande, \code{One(c,0)}. (c) La hauteur est 2 :
  toute chaîne strictement croissante a au plus trois éléments, donc la suite
  des joins est stationnaire ; le join vérifie les deux conditions de la
  \cref{def:widening}.
\end{solution}

\begin{exercice}{Calculs d'intervalles}{intervalles}
  Calculer $[1,3] \join [2,9]$, $[1,3] \meet [2,9]$, $[0,4] \join \Bot$,
  $[1,3] + [2,9]$, et dire si $[1,3] \lleq [2,9]$. Le join
  $[1,3] \join [7,9]$ est-il exact au sens de l'\cref{ex:gamma} ?
\end{exercice}
\begin{solution}
  $[1,9]$ ; $[2,3]$ ; $[0,4]$ ($\Bot$ est neutre) ; $[3,12]$ ; non, les deux
  intervalles sont incomparables (\code{partial_cmp} rend \code{None}).
  $[1,3] \join [7,9] = [1,9]$ n'est pas exact : $\gam([1,9])$ contient $5$, qui
  n'est dans aucun des deux opérandes.
\end{solution}

\begin{exercice}{Widening à seuils}{widening}
  Avec $T = \{10, 20, 100\}$, calculer successivement
  $[0,5] \widen_T [0,12]$, puis le résultat $\widen_T [0,25]$, puis le
  résultat $\widen_T [0,101]$. Calculer aussi $[0,5] \widen_T [-3,5]$.
  Combien de sauts la borne haute peut-elle faire au plus ?
\end{exercice}
\begin{solution}
  $[0,20]$ (plus petit seuil $\ge 12$) ; $[0,100]$ ; $[0,+\infty]$ (aucun seuil
  $\ge 101$). Pour la borne basse, aucun seuil n'est $\le -3$ :
  $[-\infty, 5]$. Après son premier saut, la borne haute ne prend que des
  valeurs de $T \cup \{+\infty\}$ et ne fait que croître : au plus
  $|T| + 1 = 4$ sauts. C'est l'argument de terminaison de \adr{14}.
\end{solution}

\begin{exercice}{Une worklist à la main}{worklist}
  On analyse le corps \texttt{let i = 10; while (i > 0) \{ i = i - 1; \}},
  de CFG identique à \cref{fig:interpretation-abstraite-cfg-loop}, avec
  $T = \{0, 1, 10\}$ et un widening au troisième passage par l'arc arrière,
  comme \code{analyze_cfg}. Donner l'entrée finale de l'en-tête et la valeur
  de \code{i} à la sortie. Refaire le calcul avec le widening sans seuils.
\end{exercice}
\begin{solution}
  En-tête : $[10,10]$, puis $[9,10]$ et $[8,10]$ (joins, passages 1 et 2).
  Passage 3 : le corps rend $[7,9]$ et
  $[8,10] \widen_T [7,9] = [1,10]$ (plus grand seuil $\le 7$). Passage 4 :
  la branche vraie raffine $[1,10]$ par \code{i > 0} (entier) en $[1,10]$, le
  corps rend $[0,9]$, et $[1,10] \widen_T [0,9] = [0,10]$. L'en-tête vaut
  $[0,10]$ ; la branche vraie y donne encore $[1,10]$ (inchangé, la worklist se
  vide) et la branche fausse $[0,10] \meet [-\infty,0] = [0,0]$ : \code{i}
  vaut $[0,0]$ à la sortie, la réponse exacte. Sans seuils, le passage 3 donne
  $[-\infty,10]$, qui est déjà un post-point fixe, et la sortie vaut
  $[-\infty, 0]$ : correct, mais la borne basse est perdue.
\end{solution}

\begin{exercice}{Dominance et niveau}{dominance}
  Pour \code{Conditional} (\cref{lst:interpretation-abstraite-setter}),
  dessiner le CFG du rendu, calculer $\mathrm{Dom}$ de chaque bloc, et en
  déduire ce que rendent \code{certify} et \code{may_be_skipped} pour le bloc
  de \code{setN(1)}. Retrouver le niveau observé.
\end{exercice}
\begin{solution}
  \code{lower_if} (\src{src/lowering/cfg_builder.rs}{544}{585}) scelle
  toujours un bloc \code{then} et un bloc \code{else}, même sans
  \code{else} source : quatre blocs, pas trois. L'entrée $E$ porte le
  \code{useState} et le test de \code{flag} ; $S$ est le \code{then}
  (\code{setN(1)}, puis saut) ; $X$ est l'\code{else} vide (rien à lowerer,
  saut immédiat) ; $R$ est le bloc de jonction, qui porte le \code{return} et
  n'est atteint que depuis $S$ et $X$ (jamais directement depuis $E$).
  $\mathrm{Dom}(E) = \{E\}$, $\mathrm{Dom}(S) = \{E, S\}$,
  $\mathrm{Dom}(X) = \{E, X\}$,
  $\mathrm{Dom}(R) = \{R\} \cup (\{E,S\} \cap \{E,X\}) = \{E, R\}$ : le bloc
  vide $X$ ne change rien au résultat, puisqu'il hérite tel quel du
  dominateur de $E$. $S$ ne domine pas la seule sortie $R$ : \code{certify}
  rend \code{None} (pas de \code{Certified}), \code{may_be_skipped} rend
  vrai. Sans preuve, la règle \regle{setter-in-render} s'arrête au
  \Warning{} observé.
\end{solution}

\begin{exercice}{Forte ou faible ?}{maj}
  Un objet est alloué par un littéral \texttt{const o = \{ v: 0 \}} dans le
  rendu, puis un handler fait \code{o.v = 1}. La variable \code{o} ne désigne
  qu'un site d'allocation. Pourquoi l'écriture sur le tas doit-elle
  néanmoins rester faible ?
\end{exercice}
\begin{solution}
  Le site d'allocation résume \emph{tous} les objets créés à cet endroit, un
  par rendu : un seul objet concret reçoit l'écriture, les autres gardent
  $0$. Et le handler peut ne jamais s'exécuter, ou s'exécuter sur un seul
  chemin. Remplacer $0$ par $1$ affirmerait que tous ces objets valent $1$ :
  une sous-approximation, donc un FN possible.
\end{solution}
